% 来源: 19c9fbc1-9cb2-88e6-8000-00003a43d323_3.4_调度.pdf
% 提取时间: 2026-02-28T22:51:42+08:00

3.4 调度
 我们每天都在安排各种事情，比如假期计划或者预约。路由器其实也是一样。它需要处理很多任务，
比如更新路由信息、管理查询，还有最重要的—调度数据包。数据包的调度是由路由器内部的转发处理
器实时完成的。这一章主要讲的是如何高效地调度数据包，同时让某些类型的数据包能获得比其他类型
更好的服务。




              图 1 路由器模型，其中 B3，数据包的调度
  让我们再看看路由器的工作原理（图 1）。当数据包从输入链路进来时，会通过地址查找模块找到目
标地址。地址查找会告诉路由器这个数据包应该从哪个输出链路出去，然后交换系统会把数据包送到对
应的输出端口。到了输出端口后，数据包会被放进一个FIFO队列（先进先出队列）等待发送。如果网络
拥堵，输出链路的缓冲区满了，后面新到的数据包就会被直接丢弃。这是很多路由器默认的做法，叫
做“尾部丢弃”。不过，这并不是唯一的选择：
a）我们可以根据数据包的头部信息，把它们分到不同的队列里，而不是全都塞进一个队列。
b）我们可以用一些策略来决定这些队列的优先级，比如优先级调度或者轮询调度，这样数据包的处理顺
序就不用严格按照到达的先后顺序了。
c）即使只有一个队列，也不一定要等到缓冲区满了才丢弃数据包。有时候，即使缓冲区还没满，我们也
可以提前丢弃一些数据包，这种做法可能听起来有点奇怪，但它确实有它的用途。
  为什么要搞这么复杂呢？因为通过合理的调度，我们可以为不同的数据流提供所谓的“服务质
量”（QoS），比如降低延迟或者保证带宽。要做到这一点，不仅需要调度数据包，还需要在路由器上做
一些资源预留。像RSVP和DiffServ这样的技术就是用来做这些事情的。这是因为QoS的一部分功能，比
如处理资源预留，通常由路由器的控制处理器慢慢处理。我们会聚焦在数据包调度上。

                                                   1
1.   服务质量的动因
 我们会将数据包流分配给队列尝试为流提供保证。注意，这里信息流是数据包的流，它从源到
目的经过相同的路由，并且在路径中的每个路由器或网关需要相同的服务等级。此外，流必须使用
包头中的字段来识别；这些字段通常只从传输、路由和数据链路报头中提取。流的概念是通用的，
适用于数据报网络(如 IP、OSI)和虚拟电路网络(如 X.25、ATM)。例如，在虚拟电路网络中，流可
以通过虚拟电路标识符(VCI)来标识。另一方面，在互联网上，流可以由：(1) 具有与子网A匹配的
目的地址的所有数据包识别，(2)具有与子网 B匹配的源地址的所有数据包识别，以及(3)包含特定
应用层流量的所有数据包识别。本章流识别不是重点，我们假设数据包分类算法可以有效地识别特
定流。

为什么FIFO不够，需要更高复杂性的调度算法呢? 主要原因按重要性排序如下
 ● 维持高吞吐：随着链路速度勉强跟上呈指数增长的需求，互联网上经常会出现拥塞。大多数流量基
   于TCP，它有应对拥塞的机制。然而，有了路由器的支持，可以改善TCP源对拥塞的反应，提高源
   的整体吞吐量。
 ● 流之间公平共享链路：在使用尾丢弃路由器时，客户注意到在网络备份期间，重要的Telnet和电子
   邮件连接会冻结。这是因为备份数据包占据了某个路由器输出的所有缓冲区，排除了该输出链路上
   的其他流。
 ● 为流提供QoS保证：一种更精确的公平共享形式是为流保证带宽。例如，ISP可能希望为其客户保证
   10 Mbps的带宽，作为连接客户站点的虚拟专用网络的一部分。更困难的任务是为流（如视频流）
   保证通过路由器的延迟。如果网络延迟不受限制，直播视频将无法正常工作。
可以对照一下分时操作系统(OS)(如 UNIX)。
 ● 首先，在过载时，OS必须决定要卸载哪个负载;
 ● 其次，为了公平起见，OS经常在一组相互竞争的工作中分摊时间;
 ● 最后，一些操作系统会为某些实时作业的调度提供延迟的保证，比如播放电影。



2 随机早期检测（RED）
  Random Early Detection，简称RED，是一种主动队列管理（Active Queue Management,
AQM）算法，由Sally Floyd和Van Jacobson于1993年提出，是在大多数现代路由器中实现的数
据包调度算法，即使在最高速度下，也已成为事实上的标准。简而言之，RED路由器监控平均输
出队列长度；当它超过一个阈值时，它会以一定的概率随机丢弃到达的数据包，即使可能有空间来

                                                                2
缓冲数据包。被丢弃的数据包起到了信号的作用，让源尽早减速，防止以后出现大量的丢弃数据
包。
  为了理解 RED，我们必须回顾一下互联网拥塞控制算法。图 2 的顶部显示了连接源 S 和目标
D 的网络。假设网络的链路容量为 1Mbps，文件传输速度为 1Mbps。现在假设中间的链路被一个
更快的 10Mbps 的链路所取代。数据包以 10Mbps 的速率到达第二台路由器，第二台路由器只能
以 1Mbps 的速率转发数据包；这就造成了大量的丢包，导致重传速度变慢，导致文件传输的吞吐
量非常低。
  TCP 拥塞控制的机制在发送端维护一个大小为 ( W ) 的窗口，该窗口表示在未收到确认的情况
下允许发送的数据包数量。通过调整窗口大小 ( W )，可以有效控制源端的发送速率，因为源端到
目标端的传输延迟受到发送 ( W ) 个数据包所需时间的限制。TCP 源端的窗口大小 ( W ) 初始值
为 1。在没有发生丢包的情况下，源端会以指数级增长的方式增加窗口大小，每次往返时间
（RTT）将窗口大小翻倍，直至 ( W ) 达到某个预设阈值。在此之后，窗口大小转为线性增长模
式，逐步递增。




                 图 2 依据RED实现拥塞控制
  当网络中发生数据包丢弃时（这一现象可通过接收端返回的多个具有相同序号的确认帧推断得
出），仅需重传被丢弃的数据包以填补传输序列中的“空缺”。这种机制被称为快速重传。在这种
特定场景下，发送端会感知到一定程度的网络拥塞，并将拥塞窗口大小减半（如图2所示），随后
逐步尝试增加窗口大小。然而，若多个数据包同时丢失，发送端的恢复将依赖于一个较长的、约
200毫秒的计时器超时机制。在此情况下，发送端会推断出更为严重的网络拥塞，并将拥塞窗口大
小重置为1（如图2所示）。
  例如，在采用尾部丢弃策略的路由器环境中，如图2顶部所示的示例网络可能会导致发送端的拥
塞窗口逐渐增大，直至路由器因队列溢出而丢弃部分数据包。此时，发送端会将窗口大小重置为1
并重新开始增长。尽管这种行为会导致窗口大小的周期性振荡，但由于与无拥塞控制机制的情况相
                                                   3
比，重传次数显著减少，因此发送端的平均吞吐量仍保持在较高水平。然而，如果发送端能够在每
个周期中将窗口大小降至其最大值的一半（而非完全重置为1），同时完全避免代价高昂的超时
（200毫秒），是否能够进一步优化性能？通过引入RED（随机早期检测）路由器，这种优化情景
更有可能实现。
  RED (Floyd 和 Jacobson, 1993)(图 3)的主要思想是让路由器在所有缓冲区耗尽之前尽早地检
测到拥塞，并向源发出警告。最简单的方案，称为 DECbit 方案(Ramakrishnanand Jain, 1990)，
当其平均队列大小超过阈值时，路由器将向源发送一个“标识”位。由于在当前的 IPv4 报头中没有
空间容纳这样的位，RED 路由器只是以一些小概率丢弃数据包。这使得导致拥塞的流更有可能只
丢弃一个数据包，因此可以通过更高效的快速重传而不是使传输速率变化剧烈的超时来恢复。




图 3 RED是一个早期预警系统，它隐式地通过丢包而不是像 DECbit方案那样显式地通过发送通知
随机早期检测（RED）的实现远比其表面看起来复杂。首先，我们需要利用权重为 w 的加权移动平均
值来估算输出队列的大小。假设每个到达的数据包将其观测到的瞬时队列长度视为一个样本，则平均队
列长度的计算公式为：将当前样本值乘以权重 w ，并与之前的平均队列长度乘以 (1 − w) 相加。换
言之，当 w 较小时，即使某个样本值较大，它对平均队列长度的影响也相对有限。这种设计使得平均
队列长度的变化较为平缓，需要大量的样本才能显著改变其值。这一特性是经过深思熟虑的设计选择，
旨在以往返延迟（例如 100 毫秒）的尺度上检测拥塞，而非响应短暂的瞬时拥塞。
为了简化算法实现并提高效率，可以限制 w 的取值为 2 的幂的倒数（例如 1/512 ）。虽然这种方
法在一定程度上牺牲了参数调节的灵活性，但相较于允许 w 取任意值的情况，它显著优化了计算效
率，因为乘法操作可以被高效的位移操作所替代。




                                                            4
                图 4 使用RED阈值计算丢弃概率
  然而，RED 的复杂性不止于此。丢包概率的计算依赖于如图 4 所示的分段线性函数。具体而言，当
平均队列长度低于设定的最小阈值时，丢包概率为零；当队列长度介于最小和最大阈值之间时，丢包概
率随队列长度线性增加至最大值；当队列长度超过最大阈值时，所有到达的数据包均被丢弃。同样地，
为了减少不必要的通用性（P7），可以采用合理的简化策略，例如将最大阈值设置为最小阈值的两倍，
并确保最大阈值为 2 的幂。通过这种方式，插值计算可以通过两次位移操作和一次减法高效完成。
  进一步的复杂性还体现在如何处理突发流量。如果某源端在短时间内连续发送多个数据包，现有的
RED 版本可能会连续丢弃多个数据包。为降低这种情况发生的概率，从而提升快速重传机制的成功率，
丢包概率会通过一个函数进行调整，该函数依赖于自上次丢弃以来排队的数据包数量。这种调整机制使
得丢包概率随着未丢弃数据包的数量逐渐增加，从而降低了短时间内的连续丢弃事件的发生频率。
  除此之外，还需要考虑针对不同类型流量的差异化处理需求。例如，突发流量可能需要更高的最小阈
值以避免误判。为此，思科提出了加权随机早期检测（WRED），其中阈值可以根据 IP 报头中的服务类
型（TOS）字段动态调整。此外，路由器上的随机数生成也是一个关键问题。通常的做法是利用路由器
中某些看似随机的寄存器值（例如高速时钟的低阶位）作为随机数来源。这些低阶位由于更新速度远快
于数据包到达速率，因此能提供足够的随机性。
  综上所述，尽管 RED 的核心思想看似简单，但在实际部署中却涉及诸多复杂的设计考量，特别是在
千兆级网络环境中。尽管如此，RED 凭借其良好的可行性和适应性，已成为现代路由器设计中不可或缺
的一部分。

3 近似公平丢弃（AFD）
                                                5
    有人可能会说，RED其实不是数据包调度算法，而是一种准入控制机制，因为路由器只是按照“先
来后到”的顺序（FIFO）转发所有没被RED丢弃的数据包。不过，我们后面会讨论真正的数据包调度算
法，这些算法的任务是决定多个流量或流之间如何公平地分配带宽。相比这些算法，RED及其变种（比
如思科路由器用的加权RED）计算和系统开销都比较小，但在公平分配带宽方面表现也差一些。不过，
这并不是说RED有问题：它本来就是为了拥塞控制设计的，能顺便在一定程度上实现公平带宽分配，已
经算是额外的好处了。
    因此，当“近似公平丢弃”（AFD）的技术被提出时，大家都觉得挺惊喜。AFD的计算和系统开销跟
RED差不多，但它的公平带宽分配效果却接近一些实际的数据包调度算法（比如Deficit Round Robin，
简称DRR）。虽然加权RED也可以对不同流量进行差异化丢弃，但它无法准确实现目标带宽分配（比如
按某种公平性标准分配带宽），这是它的一大局限。相比之下，AFD可以做到几乎和DRR一样精细、准
确的带宽分配。
AFD的基本原理
AFD的核心思想是通过动态调整丢弃率来实现带宽的公平分配。在无需维护所有流状态的情况下，实
现“流级别的公平丢弃策略”。以下是AFD算法的详细步骤：
1. 流量分类
必须具备“流感知”能力。大多数传统AQM算法（如RED）是包级别的，对“包”一视同仁。但在网络中，
大流（如视频或UDP洪水）和小流（如网页、短TCP连接）应区别对待，才能实现公平。AFD首先将网
络流量分为不同的类或流。分类可以基于多种因素，如源IP地址、目的IP地址、端口号、协议类型等。分
类的目的是为了能够对不同类型的流量进行独立的带宽管理。
2. 流量监控
采用哈希方法估算各流的发送速率/影响力。AFD 不维护每个流的精确状态（这在高速路由器上开销太
大），而是用哈希桶（hash buckets）来聚类流，从而估算每个流对拥塞的“贡献度”。AFD通过监控每
个类的流量到达率来估计每个类的带宽使用情况。监控通常在几个往返时间（RTT）的时间尺度上进
行，以确保估计的准确性。监控可以通过统计每个类的数据包数量或字节数来实现。
3. 丢弃率计算
拥塞时采用“选择性惩罚”策略。
 ● 如果流“贡献”大（即高速发送），那么它被丢包的概率就高。
 ● 如果流很“温和”（即低速发送），那么其包将尽量保留。




                                                       6
这种策略叫做 “惩罚大流，保护小流”（Punish heavy hitters, protect mice flows）。基于每个类的流
量到达率，AFD计算一个丢弃率，以确保每个类都能获得公平的带宽份额。丢弃率的计算类似于DRR中
的虚拟时间概念，但更为简化。具体来说，AFD使用以下步骤来计算丢弃率：
 ● 估计到达率：AFD估计每个类的流量到达率。假设在第i个类的到达率为 ri 。                ​




 ● 设置丢弃率：AFD根据估计的到达率设置丢弃率。丢弃率的设置基于一个公平性准则，例如最大-最
   小公平性（max-min fairness）。具体来说，AFD确保每个类的丢弃率 Di 使得折扣后的到达率     ​




   ri (1 − Di ) 接近其公平发送率 rif 。
     ​       ​               ​




思想很简单：时间被分成若干个长度为一个RTT的小段。路由器会测量并估计每个流量类在当前时间段
的到达率，常用方法包括，滑动窗口加权平均（EWMA）。然后根据这些数据设置下一个时间段的丢弃
率。具体来说，假设当前时间段某个流量类 i 的到达率是 ri ，那么在下一个时间段，AFD会设置一个
                                         ​




丢弃率 Di ，使得调整后的到达率 ri (1 − Di ) 等于 min(ri , rif ) , 这里 rif 是分配给流量 i 的公平
发送速率，它是通过所有流量类的到达率和链路带宽，按照最大-最小公平性原则计算。
         ​               ​       ​           ​   ​   ​




4. 动态调整
AFD动态调整丢弃率以反映网络中的实时流量变化。如果一个类的流量增加，其丢弃率也会相应增加，
从而限制该类的带宽使用。反之，如果一个类的流量减少，其丢弃率也会减少，以允许更多的带宽使
用。
5. 数据包处理
当数据包到达时，AFD根据当前的丢弃率决定是否丢弃数据包。如果一个类的丢弃率高，到达的数据包
更有可能被丢弃。这种动态调整确保了带宽的公平分配。
AFD的应用场景
AFD特别适用于需要轻量级公平带宽分配的场景，例如在资源受限的设备（如某些网络交换机或路由
器）中，或者在需要快速响应流量变化的动态环境中。
通过这种方式，AFD提供了一种在计算和系统开销方面较为轻量级的公平带宽分配机制，适用于需要快
速响应和低开销的网络环境。
    AFD的优势在于，它可以诱导出几乎和DRR一样精确的目标带宽分配，同时系统开销比DRR低得
多。比如，AFD不需要像DRR那样为每个流量类或每条流单独分配缓冲区，而DRR必须这么做（我们在
第7.2节会详细解释）。此外，AFD和DRR在处理传入数据包时的计算复杂度都是O(1)，也就是说它们的
计算效率相当。综合来看，AFD的性价比更高。目前，AFD和DRR都在思科的交换机和路由器产品中得
到了应用。

                                                                     7
4 令牌桶限速
  在讨论随机早期检测（RED）机制时，通常假设所有数据包共享一个统一的输出队列，并在队列入口
决定是否丢弃数据包。然而，一个问题值得思考：如果不区分不同流，能否为共享同一队列的流提供带
宽保障？例如，许多客户要求限制流量的速率。更具体地说，ISP可能希望限制其网络中的新闻流量不超
过1 Mbps。第二个例子是，UDP流量从路由器流向缓慢的远程线路。由于UDP源目前不对拥塞做出反
应，通过让管理员限制UDP流量小于远程线路速度，可以避免下游拥塞。幸运的是，这些带宽限制的例
子可以通过一种称为令牌桶监管的技术轻松实现，该技术仅使用每个流的单个队列和一个计数器。
  令牌桶限速（Token Bucket Policing）是一种经典的流量监管算法，用于控制数据流的传输速率，
确保其符合预定义的带宽规范。该算法通过维护一个虚拟的令牌桶来实现流量管控，其核心思想是：系
统以恒定的速率生成令牌并存入桶中，每个令牌代表一定量的数据传输权限。当数据包到达时，必须消
耗相应数量的令牌才能被放行；若桶中令牌不足，则根据策略对数据包执行丢弃或标记操作。
  令牌桶的运行机制包含两个关键参数：令牌生成速率（R）和桶容量（B）。令牌生成速率决定了长期
允许的平均数据速率，而桶容量则限制了短时间内允许的突发流量大小。这种设计使得令牌桶算法既能
保证流量的长期平均速率不超过预定值，又能容忍合理的突发传输，从而有效适应实际网络流量的波动
特性。
   注意，与漏桶算法相比，令牌桶算法的主要优势在于其对突发流量的适应性。漏桶算法强制平滑输出
流量，可能造成不必要的延迟；而令牌桶算法通过积累的令牌允许合规的突发传输，更符合许多实际应
用场景的需求。举个例子，我们可以限制一个流长期的平均速率为 100 Kbps，但允许它以最高 4 KB 的
速度短暂发送数据。因为大多数应用都会有突发流量，所以允许一定程度的突发是有好处的。而且，这
种方法对下游节点也有帮助，因为它可以减少突发流量引发的短期拥堵和丢包问题。其实现原理如图 5
所示。




                                                       8
                   图 5 令牌桶概念图
 你可以这样理解：每个流都有一个“桶”，这个桶会以指定的平均速率 R 每秒填充“令牌”。不过，桶的
容量是有限的，最多只能装 B 个令牌。如果桶满了，新来的令牌就会被丢弃。当数据包到达时，只有桶
里的令牌数量足够支付数据包的大小（以比特为单位），数据包才能被发送出去。如果令牌不够，数据
包就只能排队等待，直到有足够的令牌。因为桶最多能装 B 个令牌，所以突发流量会被限制在最多 B 比
特，之后流量会以更稳定的速率 R 每秒发送。这可以通过每个流使用一个计数器和计时器来轻松实现
——计时器用来增加计数器的值，而计数器的值永远不会超过 B。当数据包发送出去时，计数器的值会
相应减少。
 令牌桶整形有一个缺点：它需要为每个流单独设置一个队列。因为有些流可能会暂时用完令牌，必须
等待新的令牌生成，而其他流可能已经攒够了令牌。
5 多个出站队列和优先级
我们现在来看看多队列可能会用到的调度规则：
首先，得注意一点：现在我们需要根据数据包的头部信息来决定它应该被放到哪个出站队列里。这个操
作可以通过复杂的数据包分类技术来完成，也可以用更简单的方法，比如检查IP头部中的TOS字段。
其次，还要记住的是，我们仍然可以在把数据包放入正确的出站队列之前把它丢弃，这样可以实现像
RED（随机早期检测）和令牌桶这样的流量管理机制。
最后，我们遇到了一个新问题：如果有多个队列同时需要发送数据包，我们得决定下一个该处理哪个队
列，以及什么时候处理。

                                                9
图 6 单个出站队列(左)与多个出站队列(右)。基于丢弃的规则(如 RED 和警务)可以用单个队列实
现，但其他可能性，如轮询和优先级，也可以用多个队列实现。
作为多队列调度原则的最简单的例子，考虑严格优先级。例如，假设有两个出站队列，一个用于高
级服务，另一个用于其他数据包。假设数据包根据 IP包头中的 TOS位被多路分解到这两个队列。
在严格优先级中，只要高优先级队列中有一个数据包，我们总是会在低优先级队列之前服务高优先
级队列。
这里，我们可以参考操作系统的调度算法。

7 提供带宽保证
鉴于路由器可以为部分流量设置预留，我们现在回到调度器执行这些预留的问题上来。在本节中，
我们将集中讨论带宽预留，下一节再考虑延迟预留。我们将从第7.1节中的一个比喻开始，说明这
个问题；然后在第7.2节中描述一个解决方案。

7.1 偏狭的包裹服务
为了说明这个问题，让我们考虑一个假设的包裹服务，称为偏狭的包裹服务，如图7所示。两位顾
客，琼斯和史密斯，使用包裹服务通过卡车将他们的包裹送到下一个城市。
起初，所有的包裹都放在装货码头的一个队列中，如图8所示。不幸的是，装货码头的大小有限。
还发生了这样的情况：在繁忙时期，琼斯会在史密斯发货前不久发送他所有的包裹。结果是，当史
密斯的包裹在繁忙时期到达时，它们被拒绝了；史密斯被要求改天再试。
为了解决这种不公平的问题，偏狭的包裹服务决定在装货码头前使用两个队列，一个给琼斯，一个
给史密斯。在繁忙的时候，为史密斯的队列留出一些空间。队列以循环方式提供服务。不幸的是，
即使这样也没有很好地工作。
                                                10
图7 一个假设的包裹服务。




图8 一个FIFO队列用于装载包裹，不幸的是被琼斯霸占了。
这是因为邪恶的琼斯（见图9）巧妙地使用了始终比史密斯的大包裹。由于琼斯的两个大包裹可以
包含史密斯的七个包裹，结果是在繁忙时期琼斯可以得到史密斯3.5倍的服务。因此，史密斯更高
兴了，但他仍然不满意。




 FIGURE 9 Two queues and round-robin make Smith happier . . . but not completely happy.


                                                                                          11
偏狭的包裹服务短暂尝试过的另一个想法实际上是切割包裹成片，比如单位立方体，这些切片需要
标准时间来处理。然后，公司可以一次处理每个客户的切片。他们称这种方法为切片循环。当最初
的实地试验引起顾客的强烈不满时，偏狭的包裹服务决定他们不能物理地将包裹切成片。然而，他
们意识到他们可以计算出一个包裹在假想的切片系统中离开的时间。然后，他们可以按照包裹在假
想系统中应该离开的顺序来处理包裹。这样的系统对于任何组合的包裹（哦，包裹）大小来说确实
是公平的。
不幸的是，模拟假想系统就像实时执行离散事件模拟。至少，这需要记录每个队列头部包裹的出发
时间戳，并选择最早的时间戳来处理下一个；因此，这种选择所需的时间（使用优先级队列）与队
列数量的对数成正比。每当发送一个包裹时都必须这样做。
更糟糕的是，当一个新的队列变得活跃时，潜在的所有时间戳都必须改变。这在图10中展示。琼斯
有一个在他队列头部的包裹，预计在时间12离开；史密斯有一个预计在时间8离开的包裹。现在，
想象布朗引入了一个包裹。由于布朗的包裹在假想的切片系统中必须扫描一次，而其他三个切片扫
描一次，琼斯和史密斯的速度从每两片一片下降到每三片一片。这可能意味着布朗的到来会导致所
有时间戳的更新，这一操作的复杂性与流的数量成线性关系。

7.2 赤字轮询
所有这些关于包裹服务的事情是怎么回事？显然，包裹对应于数据包，包裹办公室对应于路由器，
装货码头对应于出站链路。更重要的是，看似戏谑的切片循环对应于一个开创性的想法，称为位循
环或DKS（Demers, Keshav, 和 Shenker）方案（Demers等人，1989）。模拟的位循环提供了可
证明的公平带宽分配和一些非常严格的延迟界限；不幸的是，它在千兆速度下难以实现。Stiliadis
和Varma（1996b）的论文提出了对位循环的一个重大改进，展示了如何将DKS方案的对数开销降
低到排序的对数开销。排序可以在高速下使用硬件多路堆栈完成；然而，对于带宽保证来说，赤字
轮询（DRR）仍然比DKS更复杂。




图10 布朗的到来导致琼斯和史密斯的时间戳发生变化。一般来说，当一个新的流变得活跃时，开
销与流的数量成线性关系。
                                                          12
现在，虽然位循环提供了带宽保证和延迟界限，但我们首先要观察的是，许多应用程序可以从仅仅
带宽保证中受益。因此，一个有趣的问题是，是否存在一个更简单的算法，可以仅仅提供带宽保
证。当然，我们正在放松系统要求，为更高效的实现铺平道路，正如P3所建议的。
如果我们只对带宽保证感兴趣，并希望有一个常数时间的算法，一个自然的出发点是轮询。所以我
们问自己：我们能否保留轮询的效率，同时添加一点状态来纠正图9中例子的不公平性？
一个银行类比激发了解决方案。每个流都被给予一个量子，这就像一个定期工资，每个轮询周期都
会被记入流的银行账户。与大多数银行账户一样，一个流不能花费（即发送相应大小的数据包）超
过其账户中的金额；该算法不允许银行账户透支。然而，自然地，余额会留在账户中，以便在下一
个时期可能花费。因此，任何可能的轮询中的不公平都会在随后的轮询中得到补偿，从而实现长期
的公平。
更准确地说，对于每个流i，该算法保持一个量子大小Qi和一个赤字计数器Di。分配给流的量子大
小越大，它获得的带宽份额就越大。在每个轮询扫描中，该算法将为流i服务的尽可能多的数据
包，其大小小于Qi+Di。如果流i的队列中仍有数据包，该算法会将“赤字”或剩余部分存储在Di中，
以便下次有机会时使用。很容易证明，该算法在长期内对任何数据包大小组合都是公平的，并且实
现起来只需要比轮询多几个指令。




图11 赤字轮询：开始时，所有赤字变量都被初始化为零。轮询指针指向活跃列表的顶部。当第一个队列
被服务时，量子值500被加到赤字值上。服务完队列后，剩余的量留在赤字变量中。


                                              13
图12
赤字轮询（2）：在发送了一个大小为200字节的数据包后，F1的队列中还剩下300字节的配额。然而它
无法在当前轮次使用这些配额，因为队列中的下一个数据包大小为750字节。因此，这300字节将会被顺
延到下一轮，在下一轮它可以发送总计大小为300字节（上一轮的剩余配额）加上500字节（当前轮次的
配额）的数据包。
考虑图11和图12中说明的例子。我们假设所有流的量子大小为500，有四个流。在图11中，轮询指
针指向F1的队列；该算法将量子大小500加到F1的赤字计数器上，现在为500。因此，F1有足够的
资金发送其队列头部大小为200的数据包，但不能发送第二个大小为750的数据包。因此，剩余的
300被留在F1的赤字账户中，算法跳到F2，留下如图12所示的画面。
因此在第二轮中，该算法将发送F2队列头部的数据包（留下0的赤字），F3队列头部的数据包（留
下400的赤字），以及F4队列头部的数据包（留下320的赤字）。然后它返回到F1的队列。F1的赤
字计数器现在增加到800；这反映了过去账户余额300加上新的存款500。然后该算法发送大小为
750的数据包和大小为20的数据包。假设没有更多的数据包到达F1的队列，如图12所示。因此，由
于F1队列为空，该算法跳到F2。
奇怪的是，当跳到F2时，该算法并没有留下F1队列中800-750-20=30的赤字。相反，它将F1的赤
字计数器归零。因此，赤字计数器是一个有点奇怪的银行账户，除非账户持有人能证明“需
要”（以非空队列的形式），否则会被清零。也许这类似于福利账户。

7.3 赤字轮询的实现和扩展
                                                 14
如前所述，DRR有一个主要的实现问题。该算法可能会访问许多没有数据包要发送的队列。如果
可能的队列数量远大于活动队列的数量，这将是非常浪费的。然而，有一个简单的方法可以避免空
闲跳过不活动的队列，那就是为速度添加冗余状态（P12）。
更准确地说，该算法维护一个辅助队列ActiveList，它是包含至少一个数据包的队列索引的列表。
在示例中，F1在它的最后一个数据包被服务后被从ActiveList中移除。如果F1的数据包队列不为
空，该算法会将F1放在ActiveList的尾部，并跟踪任何未使用的赤字。请注意，这防止了一个流在
空闲时在其账户中增加量子。
请注意，DRR在流之间按比例共享带宽。例如，假设有三个流F1、F2和F3，它们的量子大小分别
为2、2和3，都有预留。那么，如果这三个流都处于活动状态，F2应该获得输出链路带宽的
2/(2+2+3)=2/7。如果，例如，F3处于空闲状态，那么F2保证获得输出链路带宽的
2/(2+2)=1/2。在任何情况下，一个流都被保证在它处于活动状态的期间获得最小带宽，该带宽与
它的量子大小与所有预留的量子大小总和的比例成正比。
DRR的效率如何？只要每个流的量子大于最大尺寸的数据包，出队一个数据包的成本就是固定数
量的指令，因为只要访问一个队列，就会发送一个数据包。例如，如果一个流的量子大小为1，该
算法将不得不访问一个队列100次来发送一个大小为100的数据包。因此，如果最大数据包大小为
1500，流F1要接收的带宽是流F2的两倍，我们可以安排F1的量子为3000，F2的量子为1500。同
样，就我们的原则而言，我们注意到避免量子设置的通用性（P7）允许更高效的实现。
DRR的扩展
我们现在考虑DRR的两个扩展：分层DRR和带有单个优先级队列的DRR。
分层赤字轮询
在所谓的基于类的排队（CBQ）方案（Floyd和Jacobson, 1995）中引入了一个有趣的带宽共享模
型。这个想法是指定一个用户层次结构，可以共享一个输出链路。例如，一个跨大西洋的链路可以
由两个组织按其支付链路费用的比例共享。考虑两个组织A和B，它们分别支付链路成本的70%和
30%，因此希望在带宽中拥有相应的保证份额。然而，在组织A内部有两种主要的流量类型：Web
用户和其他用户。组织A希望限制Web流量，当A的其他流量存在时，只获得A的流量份额的
40%。同样，B希望视频流量在B的其他流量存在时不超过总流量的50%（图13）。




                                                   15
               图13 基于类别的带宽共享队列规范示例。
假设在某一时刻，组织A的流量只有Web流量，而组织B既有视频流量也有其他流量。那么A的
Web流量应该获得A的全部带宽份额（比如说，1 Mbps链路中的0.7 Mbps）；B的视频流量应该获
得剩余份额的50%，即0.15 Mbps。如果A的其他流量上线，那么A的Web流量的份额应该下降到
0.7x0.4=0.28 Mbps。CBQ可以使用每个节点的分层DRR调度器轻松实现。例如，我们将使用
DRR调度器来划分A和B之间的流量。当访问A的队列时，我们为A的流量运行DRR调度器，然后访
问Web队列和其他流量队列，并按其量子的比例为它们提供服务。
带有优先级的赤字轮询
思科系统公司实施的一个简单想法（称为修改后的DRR，或MDRR）是将DRR与优先级结合起
来，以允许IP语音的最小延迟。这个想法如图14所示，允许路由器最多有八个流队列。数据包是根
据IP TOS字段中的位（称为IP优先级位）放入队列中的。然而，队列1是一个特殊的队列，通常保
留给IP语音。有两种模式：在第一种模式下，队列1被赋予严格优先于其他队列的优先级。因此，
在图中，我们会先服务队列1的所有三个数据包，然后再在队列2和3之间交替。另一方面，在交替
优先级模式下，队列1的访问与对剩余队列的DRR扫描交替进行。因此，在这种模式下，我们会先
服务队列1，然后是队列2，然后是队列1，然后是队列3，依此类推。




                FIGURE 14 Cisco’s MDRR scheme.

8 提供延迟保证的调度器
到目前为止，我们只考虑了在多个队列之间提供带宽保证的调度器。我们的例外是MDRR，这是一
个临时解决方案。我们现在考虑提供延迟界限的问题。这种情况类似于几位厨师共用一个烤箱，如
图15所示。冷冻食品厨师（类似于FTP流量）更关心吞吐量而不是延迟；普通厨师（类似于SSH流
量）关心延迟，但对于快餐厨师（类似于语音或实时视频流量）来说，小延迟对业务至关重要。

                                                 16
FIGURE 15 Three types of chefs sharing an oven, among whom only the fast-food chef needs
bounded delay.
在实践中，大多数路由器使用DRR等算法实现某种形式的吞吐量共享。然而，几乎没有商业路由
器实现保证延迟界限的调度器。结果是，实时视频（例如远程手术）可能并不总是工作得很好。这
对于商业用途可能是不可接受的。解决这个问题的一种方法是拥有高度未充分利用的链路，并采用
MDRR等临时方案。如果带宽变得充足，这可能会奏效。然而，流量确实会增加以补偿增加的带
宽；看看Netflix导致的流量激增就知道了。
理论上，我们之前提到的模拟位循环算法（Demers等人，1989）保证了隔离和延迟界限。因此，
它被用作IntServ提案的基础，作为一种可以集成视频、语音和数据的调度器。然而，位循环，也
称为广义处理器共享，接下来将详细描述，是不现实的，因为它需要在每次服务一位后进行（相
对）昂贵的上下文切换。
WFQ，也在Parekh和Gallager（1993）中被称为分组化广义处理器共享（PGPS），提供了对
GPS的密切近似。WFQ没有昂贵的上下文切换问题，因此是现实的，因为它的调度是按数据包进
行的，而不是按位进行的。然而，正如第10节将展示的那样，WFQ微妙地依赖于GPS，如下所
述：WFQ的（调度）策略声明包含了数据包的GPS虚拟完成时间的引用，因此其实现需要跟踪
GPS时钟；这两个概念都将在第9.4节中定义。
几十年来，人们普遍认为GPS时钟跟踪操作需要对每个数据包进行禁止的高计算复杂度O(n)
（Parekh和Gallager，1993；Bennett和Zhang，1996a；Stiliadis和Varma，1996b,a；Bennett
和Zhang，1997；Chao和Guo，2001），其中n是并发流的数量。因此，WFQ一直被认为是不切
实际的；因此，已经提出了几个WFQ近似方案（例如，Zhang，1991；Bennett和Zhang，
1996a, 1997；Cobb等人，1996；Stiliadis和Varma，1996b；Suri等人，1997），以将对GPS跟
踪的复杂度降低到每个数据包O(logn)，所有这些都付出了GPS跟踪误差的代价。然而，现在已知
GPS跟踪的最坏情况计算复杂度确实可以通过Valente（2004）提出的增强数据结构将每个数据包
限制在O(logn)。有了这个，WFQ的计算复杂度也可以被限制在每个数据包O(logn)，使得WFQ可
以说与WFQ近似方案一样计算效率高。
接下来的几节（第9节至第14节）开始更严格地解释为什么在提供延迟界限的同时还要保证公平性
是一个挑战。虽然没有这些初步内容，很难欣赏Quick Fair Queuing（第14节）背后的创新，但不
                                                                                           17
太关心理论基础的实现者可能希望直接跳到第14节。

9 广义处理器共享
提供延迟和公平性保证的完美服务策略是GPS（Keshav等人，1990；Parekh和Gallager，
1993），我们之前将其描述为“模拟位循环”算法。在GPS调度器中，所有积压的流都以加权的方
式同时被服务，如下所述。在速率为r的链路上，由GPS调度器服务的每个流Fi在每一时刻t以速率
ri=r*φi/(Σj∈B(t)φj)被服务，其中B(t)是时刻t积压的流的集合。

9.1 一个简单的例子
使用下面的简单例子有助于解释GPS调度器及相关概念。我们考虑一个数据包调度实例，其中有
三个流，即F1、F2和F3，权重相等为1。假设它们共享一个服务速率为1比特每秒的链路。流i的第j
个数据包表示为pi,j。它的到达时间和长度（以比特为单位）分别表示为 ai,j 和 li,j 。设
a1,1 = a1,1 = a3,1 = 0, a1,2 = a2,2 = a3,2 = 3 ，以及 a3,3 = 4 。设
                                                                                     ​    ​




l1,1 = 7, l1,2 = 1, l2,1 = 5, l2,2 = 4, l3,1 = 3, l3,2 = 1 ，以及 l3,3 = 6 。
       ​       ​               ​         ​       ​       ​                ​




   ​               ​               ​         ​       ​        ​                ​




                               FIGURE 16 GPS graph of a simple packet arrival instance.
这个实例的服务计划，表示为GPS图，如图16所示。每个矩形对应一个数据包，每一行（矩形）
对应一个流。每行的高度对应于相应流的权重。由于所有流的权重都是1，所以所有行的高度都相
同。每个数据包的长度等于面积，在这种情况下（因为所有矩形的高度都相同）也等于矩形的长
度。例如，如图16所示，第3行包含三个长度分别为3、1和6的矩形。它们对应于F3中的三个数据
包 p3,1 、p3,2 , p3,3 。
           ​           ​   ​




在GPS下，所有三个流在时间0开始接收服务，每个流的速率为1/3比特每秒，因为所有三个流的
第一批数据包都在时间0到达。GPS图的x轴对应于虚拟时间，定义为GPS为积压流提供的服务比
特数。例如，在图16中，虚拟时间7（比特）时，所有三个流都已接收到7比特的服务。显然，对
                                                                                              18
应的实际时间为21（秒），因为在虚拟时间间隔[0,7]（比特）内，假设链路速率为r=1比特每秒，
三个流共接收到21（比特）的服务。

9.2 GPS虚拟开始和完成时间的递归定义
尽管GPS不能直接用作现实的调度器，如前所述，但与GPS相关的一些概念在指定、推理和实现
现实调度器（如WFQ、WF2Q（Bennett和Zhang，1996a）和WFQ+（Bennett和Zhang，
1997））时非常重要。其中一个概念是数据包pi,k的GPS虚拟完成时间，定义为该数据包在GPS下
完成服务的虚拟时间，表示为fi,k。例如，在图16中，数据包p3,2的GPS虚拟完成时间是4（比
特）。同样，数据包pi,k的GPS虚拟开始时间定义为该数据包在GPS下开始接收服务的虚拟时间，
表示为si,k。例如，p3,2的GPS虚拟开始时间是3。一般来说，任何数据包的虚拟开始时间和虚拟
完成时间都可以通过以下简单的递归公式计算：
                                           (1)
                  ​ fi,k = si,k + li,k /φi
                       ​   ​   ​            (2)
这里，V(ai,k)是对应于ai,k的实际到达时间的虚拟时间。例如，当数据包p3,2在时间a3,2=3到达
时，相应的虚拟时间V(3)=1/3，因为在实际时间间隔[0,1]内，有三个积压流共享GPS服务，每个流
接收1/3比特的服务。如(1)所述，数据包pi,k的虚拟开始时间定义为fi,k-1（其前一个数据包pi,k-1
的虚拟完成时间）和V(ai,k)中的较大者。这是因为，根据GPS策略，数据包pi,k在GPS下不能开始
接收服务，直到同一流的前一个数据包pi,k-1在GPS下完成服务。由于在每个GPS轮次中流Fi接收
φi比特的服务，GPS虚拟完成时间比GPS虚拟开始时间晚li,k/φi比特，这正是公式(2)。公式(1)和
(2)意味着以下事实。
事实1. 数据包的虚拟开始时间和虚拟完成时间在数据包到达时就已经确定。特别是，它们既不依赖
于未来的数据包到达，也不会随着未来数据包的到达而改变。
例如，当数据包p3,3在真实时间4（秒）到达时，它的虚拟开始时间s3,3和虚拟完成时间f3,3已知
分别为4（比特）和10，并且即使在真实时间4（秒）之后的某个时间点（比如真实时间5）有新的
流共享链路带宽，它们也保持不变。由于GPS图中的x轴对应于虚拟时间，另一种说法是，GPS图
中现有数据包（矩形）的开始和结束位置不会随着未来数据包的到达而改变。
相比之下，数据包的预期实际完成时间可能会随着未来数据包的到达而改变；这一点在第7.1节的
偏狭包裹服务示例中提到过。因此，确切地说，直到“完成”事件实际发生之前，数据包的实际完
成时间并不确定。例如，在时间0+时，数据包p1,1的预期实际完成时间是15，因为此时队列中只有
三个数据包p1,1、p2,1和p3,1，它们的长度分别为7、5和3比特，GPS调度器需要15秒来完成对p1,1
（三个中最长的）的服务。然而，在时间3+时，由于p1,2、p2,2和p3,2在真实时间3到达，其预期
实际完成时间变为18，而在时间4时，由于p3,3在真实时间4到达，其预期实际完成时间变为21。
                                                      19
此外，我们不能确定p1,1的实际完成时间确实是21，直到它的虚拟完成时间7（从时间0+开始已
知）在真实时间21到来。更一般地说，任何未来虚拟时间对应的实际时间在（真实和虚拟）时间到
来之前是不确定的。这使得跟踪GPS时钟的任务变得更加棘手，正如第9.4节将展示的那样。
如公式(1)所示，当数据包pi,k在真实时间ai,k到达时，我们需要获得相应的虚拟时间V(ai,k)。换句
话说，我们需要计算与真实时间ai,k相对应的虚拟时间。
这种虚拟到实际的映射在数学上是明确定义的，或者换句话说，不会受到前面提到的“直到时间到
来才最终确定”的问题的影响，因为这里的映射是在数据包已经到达（在时间ai,k）之后进行的。
然而，这种映射所涉及的计算，即上述GPS时钟跟踪任务的主体，并不是一件容易的事，我们将
在第9.4节详细阐述。
虽然在图16中这种计算似乎很简单，但当调度实例不那么“友好”时，情况就不是这样了。在下一
节中，我们将介绍一个不比图16中的实例大多少，但稍微复杂一点的实例。即使在这个相当简单的
实例中，GPS时钟跟踪也变得更加复杂。

9.3 另一个数据包到达场景
上述工作中提出的一个场景用于展示一种算法，该算法在最坏情况下以O(logn)的计算复杂度跟踪
GPS时钟。它与之前的简单场景在两个重要方面有所不同。首先，所有流的权重不相同。其次，
与前一个实例中GPS图中的数据包“背靠背”且没有“间隙”不同，这个场景中有一个“空洞”。现
在，当一个新的流在真实时间t到达时，确定相应的虚拟时间V(t)不再是一件简单的事情。
更准确地说，考虑一个数据包调度场景，其中有三个流，即F1、F2和F3，权重分别为1、1和2。它
们再次共享一个速率为r=1比特每秒的链路。术语pi,j、ai,j和li,j的定义与之前相同。数据包到达实
例包括p1,1，其ai,1=0且li,1=20, p2,1，其ai,1=11且li,1=10, p3,1，其ai,1=23且li,1=10，最后是p2,2，
其ai,2=39且li,2=10。




                 FIGURE 17 GPS graph of a more complicated instance.
                                                                            20
这个实例的GPS图如图17所示。第一个数据包p1,1的GPS虚拟开始和完成时间很容易计算：s1,1=0
且f1,1=20。后续数据包的时间需要一些推理。对于第二个数据包p2,1，其虚拟开始时间s2,1是11，
因为在真实时间间隔[0,11]内，F1是唯一活跃的流，在此期间获得了11比特的服务。其虚拟完成时
间f2,1是s2,1+li,1/φ2=21。当数据包p3,1在t=23到达时，相应的虚拟时间V(23)是17，因为在真实
时间间隔[11,23]内，F1和F2都处于活跃状态，每个流在GPS下接收6比特的服务。因此，
s3,1=17，由于F3的权重为2（因此数据包p3,1的高度是其他数据包的两倍），f3,1=s3,1+li,1/
φ3=17+10/2=22。
通过这个例子的解释，我们现在准备准确地说明虚拟时间的演变。设ti，i=0,1,2,…，是过去（即早
于当前时间t）发生的事件的真实时间，这些事件导致了积压流集合的变化；请注意，我们在这里
强调“过去”这个词，因为在任何真实时间t=η时，函数V(t)的图像仅针对真实时间区间[0,η]最终确
定。有两种类型的此类事件。第一种类型是，由于来自该流的新数据包到达（从该流），之前未积
压的流变为积压。例如，在图17中，当数据包p3,1到达时会发生这样的事件。
然而，在图16中，数据包p1,2在真实时间3到达并不会导致此类事件，因为在GPS图中p1,1和p1,2
是“背靠背”的（即之间没有间隙），积压流的总权重保持不变。第二种类型的事件是，之前积压
的流变为未积压。例如，在图17中，当数据包p1,1在真实时间35下在GPS下完成服务时会发生这样
的事件。
通过定义ti，很明显t0<t1<t2<…。GPS虚拟时间V(t)，作为真实时间t的函数，是根据(3)和(4)计算
的，如下所示。
 V (t0 ) = 0 , (3)
     ​




                                         , 0 ≤ τ ≤ ti − ti−1 , i = 1, 2, 3, … (4)
                                 τ
 V (ti−1 + τ ) = V (ti−1 ) +
                             Σj∈B(t ) ϕj
         ​           ​                           ​   ​   ​




                               i−1
                                         ​   ​

                                     ​




我们现在通过“归纳法”证明(4)中对V(·)的定义与我们刚才陈述的虚拟时间是一致的。假设V(ti-1)
的值等于真实时间ti-1对应的虚拟时间。那么，在区间[ti-1,ti]内，积压流的总权重是Σj∈B(ti-
1)φj。由于ti-1+τ∈[ti-1,ti]对于任何0≤τ≤ti-ti-1，我们通过GPS策略的定义（作为位循环）得到
(4)。

9.4 跟踪GPS时钟
当p2,2在真实时间t=39到达时，我们需要计算相应的虚拟时间V(39)。在图17中，我们“显示”这个虚拟时
间在f2,1=21“之后”但在f3,1=22“之前”。在实践中，这样的“之前和之后”关系必须在能够确定V(39)的值
之前确定。确定这种“之前或之后”关系的一种方法是主动计算并维护每个积压流中最后一个数据包的预
期实际完成时间，如下例所示。


                                                                               21
示例。在图17中，在真实时间t=23+时，与虚拟完成时间f1,1=20、f2,1=21和f3,1=22相对应的预期实际完
成时间分别是V^(-1)(20)=23+(20-17)*4=35, V^(-1)(21)=35+(21-20)*3=38，以及V^(-1)(22)=38+(22-
21)*2=40，分别。注意乘法项4、3和2是在真实时间间隔[17,20]、[20,21]和[21,22]内积压流的总权重。
由于a2,2=39介于V^(-1)(21)=38和V^(-1)(22)=40之间，我们知道V(39)=21+(39-38)/2=21.5（因为
V^(-1)(21)=38，且在真实时间间隔[38,39]内的总权重是2）。
这个示例还突出了之前提到的一个重要事实。也就是说，跟踪GPS时钟包括两种相互交织的任务类型。
第一种是评估V^(-1)(v)（即给定虚拟时间计算实际时间），第二种是评估V(t)（即给定过去真实时间计算
虚拟时间）。每种类型都依赖于另一种。例如，在图17所示的实例中，为了评估V(39)，我们需要评估
V^(-1)(21)。但要获得这个虚拟完成时间21（对于数据包p2,1），我们需要评估V(11)=11（当数据包p2,1在
真实时间t=11到达时）。
由于这种相互依赖性，尽管只有GPS虚拟完成时间（对于数据包）在WFQ和WF2Q的政策声明中使用，
正如第10节和第11节将分别展示的那样，也许讽刺的是，实现这两个政策的关键困难来自于在计算预期的
未来实际时间（假设在此未来时间之前不会有新流到达）时出现的困境，对应于未来的虚拟（完成）时
间，这将在下面解释。
虽然这两种类型的GPS时钟跟踪任务都可以通过主动计算并维护预期的实际完成时间来完成，如前面的
示例所示，但这样做的高昂代价是最坏情况下每个数据包O(n)的计算复杂度。这是因为，如果（新流的
第一个数据包）到达，所有预先计算的实际完成时间都将改变并需要重新计算，而且可能有O(n)个。事
实上，主动计算实际完成时间在Parekh和Gallager（1993）中被建议作为WFQ算法的一部分，用于跟踪
GPS时钟。大多数后来的论文引用了Parekh和Gallager（1993）中提到的每个数据包O(n)的复杂度，用
于执行WFQ。
然而，预先计算所有（预期的）实际完成时间可以说是“自找麻烦”。相反，似乎我们只需要预先计算其
中最早的一个来跟踪GPS时钟（P2b）。例如，在前面的示例中，当p3,1在真实时间23到达时，我们只
需要预先计算V^(-1)(20)=35，即前面描述的三个GPS虚拟完成时间（f1,1=20、f2,1=21和f3,1=22）中最
早的一个对应的实际时间。然后，直到真实时间t=35，调度器才会唤醒，并设置下一个预期的GPS实际
完成时间（即V^(-1)(21)=38）。我们通过这种方式展示了基于定时器的实现，调整总权重并安排下一个
唤醒时间，称为清理，因为每次这样的操作都对应于与前一个积压流离开相关的必要维护。
一旦调度器在真实时间t=35再次唤醒，清理操作需要将积压流的总权重从4调整为3（因为F1在t=35之后
不再积压），并设置下一个唤醒时间为V^(-1)(21)=38。这一次，调度器并没有在预定的真实时间t=38醒
来。相反，在t=38时，调度器因p3,1的到达而被“粗暴唤醒”。当这种情况发生时，调度器计算V(39)=21+
(39-38)/2=21.5，总权重调整为3（因为F1再次积压），并将下一个唤醒时间改为t=39+(22-
21.5)*3=40.5（之前为t=40）。请注意，数据包p2,1和p2,2之间存在一个小的间隙。这是因为p2,1在真
实时间t=38下在GPS下完成服务，而p2,2直到真实时间t=39才到达。
通过这种理想的基于定时器的实现，每次清理操作的时间复杂度，包括计算下一个（实际）预期唤醒时
间（可能由于“粗暴唤醒”而改变）和调整积压流的总权重，都是O(1)，如前面的实例所示。此外，我们
可以在平衡优先级队列（如堆或平衡二叉树）中维护积压流最后一个数据包的GPS虚拟完成时间，以便
                                                                               22
可以在O(logn)时间内获得下一个最早的预期GPS虚拟完成时间（通过调用ExtractMin()方法）。例如，
在前面的实例中，在时间23+时，优先级队列包含三个节点，分别以虚拟完成时间20、21和22为键。因
此，理论上跟踪GPS时钟的每次计算的总计算复杂度是每个数据包O(logn)。
然而，为了让这种O(1)的清理算法在实践中发挥作用，其操作必须紧密跟随当前时间的进展，如下所述：
当当前时间是t时，所有过去事件的清理操作必须在t之前完成，否则我们无法在O(logn)的时间内计算
V(·)。然而，可能会出现大量（例如O(n)）此类事件在短时间内（见Zhao和Xu（2004）中的一个例子）
需要清理的情况。在这种情况下，我们必须在短时间内执行O(n)次清理操作，总时间复杂度为O(n)。因
此，实践中最坏情况下的时间复杂度仍然是每个数据包O(n)。
Zhao和Xu（2004）首次详细阐述了理论与实践之间的这种不一致。不久之后，Valente（2004）提出的
GPS时钟跟踪工作解决了这种不一致，提出了一种数据结构和算法，在理论和实践之间取得了良好的平
衡。通过这种新解决方案，每次清理操作的时间复杂度较高，为O(logn)（而不是O(1)），但在“粗暴唤
醒”事件中计算V(·)的最坏情况时间复杂度也被限制在O(logn)。因此，这种新解决方案允许WFQ和
WF2Q算法的最坏情况时间复杂度都被限制在每个数据包O(logn)。

10 加权公平排队
在本节和下节中，我们将区分数据包调度策略和数据包调度算法，就像我们在操作系统文献中区分策略
和机制一样。一旦我们展示了WFQ策略可以如何简洁地陈述，以及实现WFQ算法的难度，这种区别将变
得更加清晰。
我们现在描述WFQ（Parekh和Gallager，1993），上述按数据包工作的保守调度策略，它紧密逼近GPS
策略。为了陈述WFQ策略，只需指定在繁忙期间数据包的服务顺序，因为对于所有保守调度策略，给定
任何数据包到达实例，它们的繁忙期间是相同的（因此不模糊）。同样由于WFQ的保守性质，只需指定
当当前数据包传输完成时，应该选择哪个数据包进行传输，因为在一个保守策略中不允许有空闲时间
（在连续服务两个数据包之间）。有了GPS虚拟完成时间的概念，WFQ策略关于“下一个服务谁”可以简
洁地陈述如下：
在所有当前在队列中的数据包中，选择具有最早GPS虚拟完成时间的数据包进行服务。
我们强调WFQ策略的一个微妙属性。如果我们将策略陈述中的“虚拟完成时间”替换为“预期实际完成时
间”（当需要做出“下一个是谁”的决定时），策略保持不变。这是因为不难证明，对于任何一对真实时间
τ1≤τ2，它们的相应虚拟时间满足V(τ1)≤V(τ2)：因此，基于虚拟完成时间的排序与基于实际完成时间的
排序相同。然而，从实现的角度来看，这种替换后的策略陈述是不好的。这是因为（如前所述）实际时
间对应的未来虚拟时间v在v到达之前并未最终确定。相比之下，数据包的虚拟完成时间在数据包到达时
确定，并且在之后不会改变，如第9.2节中的事实1所述。
为了说明WFQ策略的工作原理，并给读者提供如何实现它的直觉，我们使用WFQ来调度图16中所示的数
据包到达实例。不难验证，得到的服务顺序是p3,1、p3,2、p2,2、p1,1、p1,2、p2,2和p3,3。
                                                      23
WFQ被认为是GPS的近似：已经证明（例如，在Parekh和Gallager，1993中），给定任何数据包到达实
例，WFQ调度下的每个数据包的实际完成时间与GPS调度下的实际完成时间之差，上限是服务最大尺寸
数据包所需的链路速率r的时间。例如，在数据包到达实例中，不难验证，对于七个数据包中的每一个，
这个差异总是上限为7/1=7，其中7是最大数据包大小，r=1是链路速率。
接下来，我们简要讨论WFQ算法，即WFQ策略的实现。WFQ算法由两部分组成。第一部分是GPS时钟的
跟踪，通过它可以确定每个数据包的GPS虚拟完成时间。第二部分是将所有当前在队列中的数据包的
GPS虚拟完成时间存储在一个平衡优先级队列数据结构中，如堆，以便可以通过调用ExtractMin()方法来
确定“下一个是谁”，该方法每次调用（数据包）的时间复杂度为O(logn)。由于第二部分是直接的，我们
只详细描述如何实现第一部分，即GPS时钟跟踪，在第12节中。

11 最坏情况公平加权公平排队
在本节中，我们描述了另一种称为最坏情况公平WFQ或简称WF2Q的调度策略（Bennett和Zhang，
1996b）。WF2Q是对WFQ的修改，以提高其公平性。我们刚刚展示了在WFQ下，任何数据包的服务完
成时间都不会比在GPS下晚太多。然而，前者可能会比后者领先相当多的时间。由于这种短期内可能出
现的较大领先，服务不同流时可能会出现显著的公平性问题，正如我们使用WF2Q论文（Bennett和
Zhang，1996b）中的一个例子所示。
在这个例子中，有11个流，分别表示为F1、F2、...、F11。其中流F1、F2、...、F10每个权重为1，流F11权
重为10。根据这个权重分配，流F1、F2、...、F10获得链路速率的50%，流F11获得另一半。流F1、
F2、...、F10每个在真实时间0有一个数据包到达，而流F11在真实时间0、0.1、0.2、...、0.9分别有10个
数据包到达。所有数据包的长度都是1比特，链路的速率是1比特每秒。这个数据包到达实例的GPS图如图
18所示。
在WFQ下，数据包的服务顺序将是p11,1、p11,2、...、p11,10、p1,1、p2,1、...、p10,1。换句话说，在WFQ
下，F11在其所有10个数据包被服务完之前，其他10个流的所有10个数据包都不会被服务。虽然这与预期
的50/50速率分配一致，但这种分配方式的公平性听起来非常不公平：F11在其50%的分配全部完成之
前，其他10个流没有任何分配。另一种看待这种不公平的方式是，数据包p11,10在GPS下的实际完成时间
是20，但在WFQ下是10，因此前者比后者领先10。




                                                                  24
                     FIGURE 18 A motivating example for WF2Q.
这种领先时间可能会增长为O(nL)，其中n（这里是11）是流的总数，L（这里是1）是最大数据包长度，
因为对于任何n，我们可以创建一个数据包到达实例，如第18图所示，其中一个流（如这里的F11）获得
50%的带宽，而其他n-1个流（如这里的F1、F2、...、F10）平等地分享剩余的50%。对于某种调度策
略，任何数据包在该策略下的实际完成时间可能比GPS下的实际完成时间提前或滞后的最大时间称为T-
WFI（Bennett和Zhang，1996b），其中T代表（实际完成）时间，WFI代表加权公平指数。使用这个定
义，我们可以说WFQ的T-WFI是O(nL)。
既然WFQ已经是GPS的近似，读者可能会想知道是否可以设计一个比WFQ更公平的数据包调度策略。令
人惊讶的是，答案是肯定的，解决方案就是WF2Q。Bennett和Zhang（1996b）证明了WF2Q具有可能
的最小（最佳）T-WFI，即O(L)（不随n增长）。WF2Q策略可以简洁地陈述如下。
在所有当前在队列中符合条件的数据包中，选择具有最早GPS虚拟完成时间的数据包进行服务。
WF2Q与WFQ略有不同，即“下一个是谁”仅从那些在数据包队列中符合条件的数据包中选择。在任何时
间t，如果一个数据包应该在t下在GPS下开始服务，则称该数据包是符合条件的。使用相同的数据包到达
实例，我们现在展示WF2Q策略在图18所示的数据包到达实例下导致了一个更公平的服务计划。
在时间0，数据包p11,1是第一个在WF2Q下接受服务的数据包，因为它在时间0符合条件（因为其GPS开
始时间是0）。当p11,1在真实时间1完成服务时，实际时间变为1。然而，与WFQ下的计划不同，数据包
p11,2不能下一个接受服务，因为在这一点上（真实时间1），相应的虚拟时间是0.05，但p11,2的GPS开
始时间是0.1（所以它还不符合条件）。因此，下一个服务的是p1,1（假设流ID用于平局决胜）。
当p1,1在真实时间2完成服务时，p11,2符合条件，因为相应的虚拟时间是0.1，因此下一个服务的是
p11,2。通过这种推理，其余数据包的服务顺序是p2,1、p11,3、p3,1、...、p11,10、p10,1。换句话说，服
务在F11和其他10个流之间交替进行。这个计划更加公平，因为可以证明（例如，在Bennett和Zhang，
1996b中），对于任何数据包，其在WFQ下的实际完成时间既不会比GPS下的实际完成时间滞后超过
L/r，也不会提前超过L/r，所需的服务最大尺寸数据包的时间。
                                                             25
12 有效GPS时钟跟踪的数据结构和算法
如前所述，GPS时钟跟踪归结为给定真实时间t作为输入，精确计算相应的虚拟时间V(t)。Valente
（2004）提出了一种称为形状数据结构的增强数据结构，以执行此计算，每次数据包的最坏情况时间复
杂度为O(logn)（P15）。为了更清楚地描述形状数据结构及其伴随算法，我们将这个问题转化为一个更
清晰的计算问题，该问题被证明是等价的。在后一个问题中，我们需要回答一个关于楼梯的查询，如图
19所示。
12.1 楼梯查询问题
在这个楼梯中，每个台阶的宽度（run）和高度（rise）的值是精确已知的，并且假设最左边垂直线的x坐
标为0。例如，图19中第一个台阶的宽度和高度用粗体突出显示。我们的计算问题是，给定任何t>0作为
输入，我们需要找到一个垂直线的x坐标，比如说v，使得0线以下和v线“在楼梯下”的面积是t，如图19所
示。一个要求是，所提出的解决方案应该具有每次查询的低最坏情况时间复杂度O(logn)，其中n是楼梯
中的台阶数。




               FIGURE 19 Data structure (before arrival).
有了这个计算问题的指定，我们预计大多数读者能够在几分钟内想出这样的解决方案。例如，对于每个
台阶边缘的垂直线，比如说x坐标为vi，我们首先预先计算（P3b）0线和v线之间的面积，并将它们存储
在一个平衡优先级队列中。然后，给定任何t，我们可以在这些面积上执行二分搜索，以在O(logn)时间内
找到相应的v。
当强加以下三个要求时，这个计算问题变得更加具有挑战性。正如我们即将解释的，不时地，可能会根
据其右边缘的x坐标“插入”一个新的台阶到楼梯中。例如，在图20中，在图19中两个先前连续的台阶之间
插入了一个新台阶。第一个要求是，所提出的解决方案在面对这样的插入时，仍然能够在O(logn)时间内
回答问题。

                                                       26
请注意，基于预计算的解决方案将不再能够在O(logn)时间内工作，因为每个预计算的面积值都会随着新
台阶的插入而改变，更新它们都需要O(n)时间。因此，需要一个不同的数据结构和算法，第二个要求
是，当插入一个台阶时，更新这个数据结构的时间复杂度也必须是O(logn)。我们还将很快展示，不时
地，楼梯的最左边的台阶需要被移除，并且这个数据结构需要更新以正确考虑移除。第三个也是最后一
个要求是，这种更新的代价也必须限制在O(logn)。




                    FIGURE 20 Data structure (after arrival).
12.2 增强数据结构
增强数据结构 A 是基于一个传统的如二叉搜索树的数据结构构建的，我们称之为基础数据结构 B 。在
B 中，通常有一组逻辑不变量 IB 与它的数据字段相关联，当 B 更新时，需要使用 B 的函数调用
（或在面向对象编程术语中的方法）来维护。例如，当要在二叉搜索树中插入一个新节点时，需要找到
                        ​




插入点（通过对树进行二叉搜索），并且可能需要向上或向下移动一些节点（在算法文献中称为旋转
（Cormen等人，2009）），以维护以下不变量：当节点以中序遍历顺序列出时，它们的搜索键值是单
调递增的。
增强数据结构 A 通常包含几个不在基础数据结构 B 中的新数据字段。这些新数据字段与它们相关的逻
辑不变量，我们称之为 IA\B ，在更新 A 时需要维护。然而，由于 B 的方法是在不知道这些新数据
字段的情况下编程的，它们通常“不尊重” IA\B ，这意味着当调用这些方法时，它们可能会破坏 IA\B
                ​




。例如，正如我们将在下一小节中展示的，当AVL树（Adelson-Velsky和Landis，1962）作为基础数据
                                    ​                            ​




结构时，插入和删除会导致旋转，这些旋转会破坏增强数据结构中与新数据字段相关的不变量。因此，
在编程增强数据结构 A 时，我们通常需要修改从 B 继承的一些方法的实现，以修复这些方法对这些不
变量造成的损害。




                                                                27
                               图21 形状与树的关系。




     FIGURE 22 The “staircase” representation of the GPS graph shown in Fig. 16.
12.3 “形状”数据结构
在本节中，我们专注于Valente（2004）提出的“形状”数据结构的设计，用于跟踪GPS时钟。我们首先
将GPS时钟跟踪问题转化为上述“二叉搜索楼梯”问题，然后描述形状数据结构如何解决后者。
GPS图的楼梯表示
上述楼梯对应于GPS图在特定真实时间 t 的某种表示。这种表示与图16-18中使用的不同之处在于三个
方面。首先，我们将积压流（在时间 t ）的顺序从上到下排序，根据每个流中最后一个数据包的GPS虚
拟完成时间。其次，我们不在任何两个“相邻”流之间留出垂直空间，因此虚拟完成时间（按升序排列）

                                                                                   28
与这些流的横向边界一起形成了楼梯的轮廓。第三，我们省略了同一流中相邻数据包的所有垂直边界，
以便每个流在图中看起来像一个“无缝的板”。这种省略在技术上是合理的，因为这些边界在GPS时钟跟
踪中没有算法意义。
我们在图22中展示了这种楼梯式的GPS图表示，楼梯（轮廓）用粗体突出显示。该图对应于图16中所示
的数据包到达实例的GPS图，在时间 t = 3 （所有七个数据包都已到达）时。流F1，其最后一个数据包
的GPS虚拟完成时间最早为8，是顶部的第一个“板”。其次是F2和F3，它们的GPS虚拟完成时间分别为9
和10。
GPS图的楼梯表示使我们能够“可视化”上述基于定时器的O(1)复杂度GPS时钟跟踪解决方案的缺陷，如
下所述。在任何GPS图表示中，当前的虚拟时间线（即虚拟时间线）从左向右移动。有缺陷的算法要求
在虚拟时间线移动到台阶的垂直边（“上升”）的右侧时，移除楼梯的最左边台阶，作为上述清理操作的
一部分。换句话说，清理操作必须“沿着陡峭的楼梯快速下行”。然而，这个要求是有问题的，因为楼梯
的一段可能是“陡峭的”，有许多O(n)个“小台阶”，对应于上述情况，即许多积压流的最后一个数据包的
GPS完成时间在一个很短的时间间隔内。
“楼梯查询处理”由形状数据结构
形状数据结构有三个主要设计目标。第一个目标是允许清理操作在需要时远远落后于虚拟时间线的进
展，以便每次清理都可以在“方便的空闲时间”（即不受虚拟时间线进展的持续压力）下完成，并且在整
个时间间隔内，清理操作的总体最坏情况时间复杂度保持在O(logn)。这个目标如果实现，将直接解
决“沿着陡峭楼梯快速下行”的缺陷。第二个目标是在“粗鲁唤醒”（在时间t）时，尽管清理操作可能远远
落后于虚拟时间线，但每次评估V(t)的最坏情况时间复杂度仍保持在O(logn)。第三个目标是允许新到达
的流（导致“粗鲁唤醒”）的第一个数据包（目前也是最后一个数据包）的GPS虚拟完成时间插入形状数
据结构，最多花费O(logn)时间。这个操作对应于上述在楼梯中插入新台阶（板）的操作。
这三个目标对应于上述“楼梯查询”计算问题的三个要求。楼梯中的情况正是GPS图在上述真实时间t时的
表示，当需要计算相应的虚拟时间V(t)时。楼梯中的每个台阶对应一个流，其高度是流的权重，其x坐标
是流中最后一个数据包的GPS虚拟完成时间。在虚拟时间线V(t)之前的流不再积压，而在真实时间t时的
流则积压。
假设这个楼梯包含m个台阶（流），其x坐标（虚拟时间）按升序排列为 v1 < v2 < ⋯ < vm ，如图
21(a)所示。我们还假设在任何新流到达之前的最后一次事件（“粗鲁唤醒”）发生在时间 t0 ，相应的虚
                                                   ​   ​       ​




拟时间 v0 = V (t0 ) 小于 v1 ，即楼梯中最早的GPS完成时间；这个假设将在我们描述形状数据结构后
                                                           ​




放宽。有了这个假设， v0 是楼梯的左边缘的x坐标，如图21(a)所示。这个假设有以下微妙的含义：每个
      ​    ​       ​




流，比如说从上数的第i个流（台阶），在 v0 和 vi 之间积压。换句话说，楼梯下的整个区域是“实心
               ​




的”（即没有间隙）。因此，楼梯下 v0 和 V (t) 之间的面积等于 t − t0 。因此，我们计算 V (t) 的
                               ​   ​




问题等同于上述“楼梯查询”问题（其中 t0 = 0 且 v0 = V (t0 ) = 0 ）。
                       ​                       ​




                           ​           ​   ​




正如前面提到的，形状数据结构是一种增强数据结构。它的基础数据结构是一个平衡搜索树，如AVL
（Adelson-Velsky和Landis，1962）或红黑树（Cormen等人，2009），有一个特殊的规定：所有键都

                                                                   29
只存储在叶节点中。众所周知，这个规定可以在不增加数据结构（任何方法的）渐近计算复杂度的情况
下实现。例如，在B+树（Elmasri和Navathe，2010）中，所有键都存储在叶节点。对于图21(a)中所示
的楼梯，树恰好包含m个叶节点，这些叶节点中包含的键值正是这m个值。更准确地说，如果我们对该树
进行中序遍历并“打印”仅叶节点的键，输出将正好是 v1 , v2 , v3 , ⋯ , vm 。                  ​   ​           ​                   ​




代替键，每个内部节点保留其子树根节点的最小和最大键值的范围信息，换句话说，是其子树中键值的
范围。这个范围信息有两个用途。首先，在没有键的情况下，内部节点依靠这些信息来执行二叉搜索。
例如，在插入一个具有特定键值的新叶节点时，需要找到插入点（通过对树进行二叉搜索），并且可能
需要向上或向下移动一些节点（在算法文献中称为旋转），以维护上述不变量。其次，这个范围的宽
度，或者更准确地说，两个键值之间的差异，用于计算增强数据结构的一个附加（相对于基础数据结
构）数据字段，称为面积。由于只有叶节点包含键，维护树的平衡的算法步骤，如各种旋转操作（例如
AVL树中的RR、RL、LR、LL旋转），需要适当地修改，如B+树中那样。
详细的数据结构设计
我们现在描述形状数据结构。我们递归地描述它，因为其基础数据结构，一个平衡二叉搜索树，最好通
过递归来解释。形状数据结构，对应于图21(a)中所示的楼梯，如图21(b)所示，其中展示了两个顶层的详
细信息，其余部分被省略。如前所述，由L0为根的（整个）树包含m个叶节点，键值为
 v1 , v2 , v3 , ⋯ , vm 。因此，根节点L0的（键）范围是 (v0 , vm ] ，其中 vk 是范围的“中间”点。这个
 vk 大致是“本应是”的中间点，如下所述：如果这棵树是一个普通的平衡二叉搜索树，其中每个节点，
  ​               ​       ​       ​                                       ​                       ​           ​




无论是内部的还是叶节点，都与一个键值相关联，那么根节点L0将与键值 vk 相关联，或者与排名顺序
      ​




中“接近”的某个值相关联（例如， vk−2 ）。
                                                                                                          ​




                                              ​




我们现在继续到树的下一个层次。其左子节点L1的范围是 (v0 , vk ] ，因此左子树包含k个叶节点，键值
为 v1 , v2 , … , vk 。其右子节点L2的范围是 (vk , vm ] ，因此右子树包含 m − k 个叶节点，键值为
                                                                              ​               ​




 vk+1 , vk+2 , … , vm 。我们不给这两个范围命名中间点，因为图21(b)中没有这样的细节。有了这个中
          ​                   ​                   ​       ​




间点字段，对二叉搜索树的二叉搜索就很直接了：将搜索键与中间点（例如L0处的 vk ）进行比较，然
              ​                           ​




后根据比较结果在左子树或右子树中继续搜索。
                                                                                                                  ​




现在，我们描述形状数据结构中的三个“增强”数据字段。为此，我们首先描述一个我们想要实现的具体
目标。根节点联合的实线封闭区域（IV）和（V）对应于根节点维护的形状信息。我们想知道它的面积是
小于还是大于 t − t0 ，即自从上次评估函数 V (⋅) （在时间 t0 ）以来经过的实际时间；“等于”的情
况是微不足道的，因此我们在这里和后续讨论中忽略它。搜索将在“小于”的情况下在左子树中进行，在
                                      ​                                               ​




其他情况下在右子树中进行。虽然自然地包括数量（IV）+（V）作为一个新的数据字段，但从数据结构
和算法设计的角度来看，这个数量在节点插入或删除时很难维护。相反，我们维护另外两个数量，从这
些数量可以轻松（即在O(1)时间内）计算出这个数量。其中一个数量是实线封闭区域（I）的面积，用原
始论文中的W表示。另一个是楼梯的高度，即流 F1 , F2 , … , Fk 的总权重，因此在图21(b)中表示为
 [1..m] 。滥用符号一点，我们将区域（IV）的面积也表示为（IV），将区域（I）的面积也表示为
                                                      ​       ​                           ​




（I），依此类推。有了这个符号滥用，这三个数量（（IV）、[1..m]和（I））之间的关系是（IV）=（
 vk − v0 ）* [1..m] - （I）。基础数据结构字段（键）范围和两个增强数据结构字段面积和高度的值，
      ​               ​




                                                                                                                      30
对于图21(b)中的树节点L0、L1和L2，如图所示。根节点的面积关联的区域是L0.area = （I）+ （II）+
（III）。这个面积用于计算区域（VI）+（V），在比较之前应该首先与 t − t0 进行比较（如果有的
话）。它的（键）范围是（ v0 , vm ]（以 vk 为“中间点”），如前所述，其高度区域的面积显然是
                                                                        ​




[1..m]。其左子节点L1的这些三个字段的值是L1.area = （I），L1.range = （ v0 , vk ]，和L1.height =
                             ​   ​               ​




                                                                                ​       ​




       k
[1..k] ∑ i 。其右子节点L2的这些值是L2.area = （II），L2.range = （ vk , vm ]，和L2.height = [(
            ​                                                               ​       ​




      i=1
                m
k+1     )..m] ∑ i 。不难检查L0、L1和L2的这三个字段的值满足以下三个方程：L0.area =
                    ​




              i=k+1
L1.area + L2.area + |L2.range| * L1.height，(5) L0.height = L1.height + L2.height， (6) L0.range =
L1.range * L2.range。 (7) 更一般地，对于任何节点及其两个子节点，三元组中的这三个字段的值满足
上述三个方程。（第一个方程对应于Valente（2004）中的“定理1”，其中“面积”被“ δW ”所替代。）
这三个方程正是增强数据结构的附加数据字段需要满足的一组逻辑不变量。当“=”（等于）被“:=”（赋
值）替换时，得到的三个赋值命令可以用于修复这些不变量可能因对基础数据结构的一些“破坏性”方法
调用（如节点插入和删除）而受到的损害。
在介绍增强数据结构（三个附加数据字段）和相关不变量时，我们选择了一个最方便且最清晰的表示方
式。然而，这种表示方式在设计描述在上面运行的“楼梯查询”算法时显得笨拙。例如，使用这种表示方
式，算法从根节点开始，检查 t − t0 < ∣L0 .range∣ ∗ L0 .height − L0 .area 。如果是这样，它继续
到左子节点L1并检查 t − t0 < ∣L1 .range∣ ∗ L0 .height − L1 .area 。如果答案是否定的（并且排除
                                         ​   ​           ​          ​




了相等的情况），那么算法必须回溯到父节点L0，然后到L1的左子节点进行下一次（二分搜索）比较，
                         ​           ​               ​       ​




这虽然不完美（但由于二分搜索的时间复杂度仍然是O(logn)）但并不错误。因此，为了便于干净的算法
设计，我们让父节点“借用”（缓存）其子节点中这三个数据字段的值，以便上述第二次比较可以在L0而
不是L1上进行。通过这种“借用”，二分搜索可以像往常一样从根节点向下遍历到叶节点，而不需要“回
溯”。




                                                                                             31
                        FIGURE 23 Data structure (after arrival).
楼梯的左边缘对应于在最近的清理操作中从形状数据结构中删除的流记录的GPS完成时间。当发生这样
的清理时，如何修复形状（增强）数据结构的逻辑不变量将在稍后讨论。到目前为止，我们假设
 v0 < v1 ，以便二分搜索（对于V(t)线）从楼梯的左边缘精确开始。在实践中，我们几乎肯定会遇到
 v0 > v1 的情况，因为正如前面解释的，形状数据结构的目的正是允许清理的时间线远远落后于V(t)线
      ​   ​




（因此函数V(·)最后一次评估的时间 t0 ）。我们现在描述 v0 > v1 的情况（我们仍然有
      ​   ​




 v1 < v2 < … < vm ），如图23所示。为了解决这种情况，我们需要对算法逻辑进行两个更改。首
                        ​                    ​   ​




先，我们用 vc = V (tc) 替换树节点范围字段中的每个 v0 出现，其中 vc 和 tc 是最近清理操作的虚
      ​   ​   ​




拟时间和真实时间，分别。例如，L0的范围现在是 (vc, vm ] 而不是 (v0 , vm ] 。其次，算法逻辑在以下
                                     ​




方面进行了扩展：算法不仅要“按需”映射t到V(t)，还要输出区域（III）+（IV）的值，如图23所示。因
                                         ​               ​   ​




此，我们的“归纳假设”是这个逻辑在时间 t0 时成立，即算法在时间 t0 时输出了 v0 = V (t0 ) 和区域
（III）的值。基于这个“归纳假设”，我们现在展示如何修改后的算法提供扩展逻辑（输出V(t)和（III）+
                            ​                        ​           ​   ​




（IV）在时间t）。回想一下，我们可以查询形状数据结构，该结构编码了这个楼梯，以找到一个虚拟时
间线v，使得楼梯下V(tc)和v之间的面积等于（I）+（II）。由于楼梯在V(t0)线的右侧是实心的（因为在
 t0 和t之间没有新的流到达），我们有（II）= t - t0。然而，（I）的面积不一定等于t0 - tc，因为楼梯
下tc和t0之间的面积不一定是实心的。相反，（I）的面积可以计算为包围矩形（( (v0 − vc) ∗ [1..m]
  ​




）减去（III），其中（III）是已知的，感谢“归纳假设”。因此，我们的算法是查询楼梯，以找到一个虚拟
                                                                 ​




时间线v，使得楼梯下vc和v之间的面积是（ v0 − vc ）* [1..m] - （III）+ t - t0。正如刚才解释的，形
状数据结构可以处理这个查询，因为它对应于第一种情况。
                                ​




这个v是我们正在寻找的V(t)的值。一旦我们得到了V(t)的值，（IV）的面积就可以计算为（V(t) - v0）*
[1..m] - (t - t0)。现在，我们已经得到了（III）+（IV）的值，所以“归纳证明”完成了。
                                                                         32
在插入和删除下的不变量维护
在本节中，我们首先描述会破坏增强数据字段不变量的两种方法调用，即流记录的插入或删除，以及如
何修复这些不变量。然后，我们描述GPS时钟跟踪操作如何触发这两种方法调用。
我们只描述插入操作，因为删除操作类似。节点的插入首先由作为平衡二叉搜索树的基础数据结构处
理。请注意，在我们的特殊情况中，所有键都存储在叶节点中，插入的新节点必须是叶节点。在平衡二
叉搜索树（如AVL或红黑树）中插入新叶节点会触发一系列旋转（例如AVL中的LL、LR、RL和RR类
型），这些旋转首先向上遍历树，然后可能向下遍历（但不一定一直到底部的另一个叶节点）。这将破
坏受旋转和从最低共同祖先到树根的路径影响的节点集合S中的增强数据字段的不变量。因此，我们需要
修复这条“破坏路径”上所有节点的这些不变量。由于在AVL或红黑树等平衡树数据结构中，S中的节点总
数和树的高度都是O(logn)，不变量修复操作的时间复杂度是O(logn)。
最后，我们描述GPS时钟跟踪操作如何触发插入和删除。插入和删除可以通过三种不同的方式发生。首
先，每次清理操作都会导致一个叶节点被删除，该叶节点对应于二叉搜索树中最左下角的叶节点，该节
点的键值是最低的。其次，一个新流在时间t到达，如图20所示，将导致一个“楼梯查询”，以计算V(t)，
即该流第一个数据包的GPS虚拟开始时间。其GPS虚拟完成时间需要插入到形状数据结构中。第三，一
个新数据包到达，比如说 pi,k+1 （当前积压流 Fi 中的第 k + 1 个数据包 pi,k+1 ）在时间t到达现有
的（即当前积压的）流 Fi ，将导致插入（叶节点，键值为 fi,k+1 ，即 pi,k+1 的GPS虚拟完成时间）
                       ​     ​                       ​




和删除（叶节点，键值为 fi,k ，即前一个数据包 pi,k 的GPS虚拟完成时间），除了计算V(t)的“楼梯查
               ​                             ​   ​




询”。请注意，在这种情况下，我们必须有V(t) < fi,k ，否则 Fi 在时间t不是积压的。
                   ​                 ​




                                 ​       ​




我们注意到，在实际操作中，WFQ调度器通常可以“摆脱”不执行任何清理操作，原因如下。要调度的链
路的利用率水平通常远低于100%。在这种情况下，链路的忙期（对于链路的组合队列）通常不是很长，
因此二叉搜索树中的叶节点数量不会大到足以开始造成麻烦（例如，超出可用内存或使树“过高”）。
然后，“抛掉”（例如，通过发出“memset”+“free”系统调用）整个增强数据结构在忙期结束时的做法将
解决问题。可以想象，在实际操作中，这种“宽松”的清理方法可以大大减少清理时间。

13 实现WFQ和WF2Q
Valente（2004）证明了仅形状数据结构就足以实现WFQ。然而，WF2Q的实现比WFQ更难，因为除了
跟踪GPS时钟和排序GPS虚拟完成时间之外，其实现还必须以某种方式维护合格数据包的集合。Valente
（2004）证明了另一种增强数据结构和算法，由Stoica和Abdel-Wahab（1995）提出，用于实现另一
种称为最早合格虚拟截止时间优先（EEVDF）的调度策略，可以用于执行每个WF2Q操作，即在最坏情
况下以O(logn)时间找到（下一个）合格的数据包，该数据包具有最低的GPS虚拟完成时间。
Valente（2004）提出的WF2Q实现包含形状数据结构和EEVDF数据结构。每个到达的数据包首先由形
状数据结构处理，以获得其虚拟开始和完成时间，如上所述，然后插入到EEVDF数据结构中，用于

                                                          33
EEVDF调度；该数据包的虚拟开始和完成时间被视为EEVDF调度的合格时间和截止时间。我们将这种数
据结构和算法总结为第18节中的练习题。

14 快速公平队列（QFQ）
回想第11节，WF2Q调度算法可以以最小的T-WFI（在某一策略下，任何数据包的实际完成时间可能比
GPS下的实际完成时间提前或滞后的最大时间，如第11节所定义）紧密逼近GPS调度器，时间复杂度为
O(logn)，使用第12节中描述的复杂数据结构（用于跟踪GPS时钟和排序GPS虚拟完成时间）和第13节中
描述的算法（用于检查资格）。在本节中，我们描述了一种称为QFQ（快速公平队列）的调度算法，它
可以实现与WF2Q几乎相同的T-WFI，但时间复杂度仅为O(1)。
1 公平轮询（FRR）算法
公平地说，QFQ并不是第一个几乎与WF2Q一样公平但时间复杂度为O(1)的调度算法。公平轮询（FRR）
（Yuan等人，2009）是第一个这样的算法。接下来，我们首先描述FRR，因为QFQ建立在FRR的基础上
并加以改进。FRR的设计基于以下见解：调度权重相同或相似的n个流在计算上要比调度权重非常不同的
流容易得多（如第11节中F11和其它10个流之间的权重差异可能相差O(n)倍）。这个见解可以从以下两个
事实中看出：（1）当所有流的权重相同或相似时，即使是轮询调度器（如DRR，时间复杂度为O(1)）的
T-WFI也不会随着n的增长而增长；（2）然而，当流的权重非常不同时，即使WFQ（时间复杂度为
O(logn)）的T-WFI也会随着n的增长而线性增长，如第11节所示。
基于这个见解，FRR的想法是将流捆绑成组，使得在每个组内，流的权重相似。分组是按照以下“指
数”方式进行的：每个流 Fk 属于组 i = logC ϕk ，其中 ϕk 是 Fk 的权重。这里C是一个控制FRR时
间复杂度和服务质量保证之间权衡的参数；较大的C导致组数较少，从而获得更好的时间复杂度，但也导
                 ​            ​           ​   ​




致更差的QoS保证（更具体地说，更差的延迟界限）。C通常设置为整数值，如2。为方便起见，我们在
下文中假设C=2。在这种情况下，每个组内的流权重最多相差2倍。分组的一个效果是组数G与n无关，或
者说相对于n是O(1)。不难证明 G = logC (ϕmax /ϕmin ) ，其中 ϕmax 和 ϕmin 分别是流的最大和最
小权重。当C=2时，在所有可想象的使用案例中，我们有G ≤ 64。
                          ​       ​   ​           ​   ​




在将流分组后，FRR采取了以下双层方法进行成本效益高的数据包调度：它对组间进行细粒度调度，对
组内流间进行粗粒度调度（P3b），如下所述。FRR将每个组（流）视为一个超级流，并使用WF2Q的一
个变体来服务这些G个超级流。这个变体，其时间复杂度为O(G log G)，比标准的WF2Q（时间复杂度为
O(log G)）计算上更昂贵，因为在这种使用情况下（FRR的WF2Q），组的权重可以动态变化（例如，当
组的组成流之一重新积压时），而在标准使用情况下，流的权重从不变化；然而，由于G很小（G ≤
64，如上所述），O(G log G)在技术上仍然是O(1)。注意，这种组间的细粒度调度（使用WF2Q的变
体）对于FRR实现良好的公平性是必要的，因为（超级流的组成流）不同超级流可以有非常不同的权
重，而在它们之间使用粗粒度调度器可能导致T-WFI随n线性增长的糟糕情况，如上所述。

                                                               34
一旦FRR选择一个组进行服务，它将根据扩展形式的DRR，按照轮询顺序选择该组所属流之一的HOL数
据包进行服务；请注意，即使是轮询调度器在对组内流进行调度时也能实现良好的公平性，因为这些流
的权重在设计上是相似的。通过这种双层方法，FRR实现了出色的整体公平性，同时具有技术上为O(1)的
低时间复杂度。
2 QFQ如何改进FRR
QFQ通过三项修改改进了FRR，将FRR的技术上为O(1)（实际上是O(G log G)）时间复杂度降低到真正的
O(1)，换句话说，消除了其对G的依赖，从而实现了QFQ。第一个修改是使用WF2Q+（Bennett和
Zhang，1996a）中提出的近似虚拟时间函数V(t)。WF2Q+与WF2Q几乎相同，只是它的V(t)只是近似跟
踪GPS时钟（而WF2Q则精确跟踪）。这种近似，加上即将描述的第三个修改，允许WF2Q+时钟以真正
的O(1)（独立于n和G）时间复杂度进行跟踪。权衡（P3b）是，其公平性保证，以T-WFI衡量，略逊于
WF2Q。
第二个修改是，它不仅根据流的权重 ϕk ，还根据其最大数据包大小 Lk 将流 Fk 分配到一个组。更准
确地说，它将流 Fk 分配到流组i，其中 i = log2 (Lk /ϕk ) （基数C=2）。Checconi等人（2013）表
                         ​                     ​       ​




明，使用WF2Q+时钟的这种修改导致QFQ的一个很好的特性，这是其真正O(1)时间复杂度的关键：系统
            ​                      ​   ​   ​




中任何两个流的HOL数据包的WF2Q+虚拟完成时间之差被一个小常数 δ 所限制，无论流量到达模式如
何。这个特性的一个较弱版本被称为全局有界时间戳（GBT）属性，在Bennett和Zhang（1996a）中提
出。
第三个修改是，为了对HOL数据包的WF2Q+虚拟完成时间进行排序（在每个组内），在每个组i中，
WF2Q+虚拟完成时间被四舍五入到某个单位值 σi 的倍数。对于每个组索引i， σi 的值被定义为 2i
（位）。通过这种 σi 整数舍入，任何HOL数据包的舍入WF2Q+虚拟完成时间只能取少量可能的值，这
                               ​                   ​




要归功于上面描述的GBT属性。因此，对任何组内舍入WF2Q+虚拟完成时间（以决定哪个流的HOL数据
包将被传输）的排序可以在O(1)时间内使用桶排序（P14）完成。在本节的其余部分，我们省略了“舍入
WF2Q+”的限定词，理解为所有提及的虚拟完成和开始时间都具有这个限定词。
对于属于组i的任何流Fk，由于其最大归一化数据包长度Lk/φk小于2σi，这是由于我们如上所述将流分配
到组的方式。这个属性反过来保证了任何数据包的虚拟完成时间不会比其虚拟开始时间晚超过2σi。对于
每个组i，定义s(i)为所有流（它们的HOL数据包）中最早的虚拟开始时间。为了实现O(1)时间复杂度，
QFQ做出了以下简化的假设：f(i)被假定为s(i) + 2σi。可以说，如果一个数据包的长度远小于其所属流的
最大数据包长度，这个假设对这个组的HOL数据包有些不公平。然而，由于组内数据包仍然按照其实际
虚拟完成时间的顺序提供服务，这种假设对整体公平性的总体影响通常很小（P3b）。
QFQ的调度策略现在可以精确地陈述如下：每当链路变为空闲时，QFQ策略在那些应该在WF2Q+时钟下
开始服务的合格组中，选择具有最早虚拟完成时间的组进行服务。直观地说，就像在WF2Q策略的实现
中一样，实现QFQ策略的调度算法需要同时进行排序（虚拟完成时间）和资格检查。

                                                               35
QFQ算法设法在这两种检查中都做到O(1)时间。关键思想是，在任何时候，将G个组划分为四个子集，根
据一个组是否准备好（R）接收服务或被阻塞（B），以及一个组是否有资格（E）或没有资格（I），并
为这四个子集维护不变量（在任何时候）。由于G ≤ 64，我们可以用G ≤ 64个不同的位置来表示这些G
个组；因此，每个子集（这些G个组的）可以用一个64位长的位图（P14）来编码，其中每个位的值表示
相应的组是否在该子集中。在这些不变量在任何时候都得到维护的情况下，QFQ算法归结为在每次调度
或维护事件中执行以下基本操作几次：在一个子集中识别具有最小索引的组。但这相当于在编码子集的
64位整数（位图）中“找到第一个”（FFO，第13.16节中提到）。由于大多数现代CPU都有一个内置的
FFO指令（P4c），该指令在一个时钟周期内完成，QFQ算法可以在几个时钟周期内处理每个调度或维护
事件。对于一个调度事件，在QFQ算法使用FFO识别出具有最小索引的组之后，它再次使用FFO在该组
内识别具有最小虚拟完成时间的流。
第一个子集是ER（合格且准备就绪）子集。每当QFQ算法选择下一个要服务的组时，它从这个子集中选
择。在任何时候t，ER子集维护两个不变量。首先，它包含“正确的”组，比如说组i*，在时间t选择（在时
间t），在这个意义上，组i*既有资格又具有所有合格组中最早的虚拟完成时间。其次，ER中所有组的虚
拟完成时间按其组索引排序，如下所述：对于任何i1, i2 ∈ ER且i1 < i2，组i1的虚拟完成时间不晚于组i2
的虚拟完成时间。有了这两个不变量，为了找到“正确的”组，QFQ算法只需要在ER的64位长位图上执行
FFO。
然而，如果没有额外的“魔法”，显然有些事情是错误的：肯定会有一个有资格的组的索引小于i*，在这种
情况下，FFO不会定位到“正确的”组i*。QFQ的巧妙之处在于阻塞所有这样的组，并将它们推到第二个子
集，称为EB（合格且被阻塞）。与ER子集相比，EB子集不包含“正确的”组i*，满足上述第二个不变量，
如Checconi等人（2013）在定理5中所证明的。另外两个子集是IB（不合格且被阻塞）和IR（不合格且
准备就绪）。Checconi等人（2013）展示了IB和IR组满足两个不变量：（1）这些组的虚拟开始时间按其
组索引排序；（2）这些组的虚拟完成时间按其组索引排序。注意（1）和（2）可以同时为真，因为我们
假设f(i) = s(i) + 2^{\sigma_i}，如上所述。
随着时间的推移，一个组可能需要从一个子集移动到另一个子集，每次这样的移动对应于一个上述维护
事件。正如Checconi等人（2013）所示，有四种类型的移动：（1）IB→IR；（2）IR→ER；（3）
IB→EB；（4）EB→ER。QFQ算法有两个显著特点：（1）在这四种移动情况下，可以维护有关这四个
子集的所有四个不变量；（2）每次这样的移动的维护成本（时间复杂度）都是真正的O(1)（因为它主要
涉及FFO操作）。正如Checconi等人（2013）所示，这两个特点都源于（1）假设f(i) = s(i) +
2^{\sigma_i}；（2）上述GBT属性；（3）分组规则i = \log_2(L_k/\phi_k)。由于其真正的O(1)时间复
杂度和优秀的QoS保证（接近WF2Q），QFQ算法已被集成到Linux内核中（QFQ源代码在Linux内核
中，2012）。

15 面向可编程数据包调度

                                                                 36
我们在第11.15节中在可编程IP查找和P4的背景下介绍了可编程交换机和路由器的概念。正如第11.15节中
提到的，P4目前不支持可编程排队规则（即数据包调度策略），但已经投入了大量研究努力来实现这一
目标。在本节中，我们介绍可编程数据包调度的概念，并描述两个代表性的研究提案。
15.1 推入式先出（PIFO）框架
在PIFO框架下，Sivaraman等人（2016a），数据包调度是可编程的，因为可以从模板中实例化几种不同
的数据包调度算法，通过选择不同的参数值。例如，配备了支持这种可编程能力的数据包调度算法的数
据中心网络，可以让网络运营商实例化最适合网络的调度算法。
PIFO框架可以支持工作和非工作守恒的数据包调度算法。然而，由于前者是本章的重点，我们在本节的
其余部分只考虑前者，并隐含地理解我们只指的是工作守恒算法。PIFO的设计基于以下观察：数据包调
度算法只需要对数据包到达做出一个决定：现有数据包（当前在数据包队列中的数据包）应该以什么顺
序被服务？对于许多算法来说，这个决定可以在数据包入队时做出，因为对于任何两个即将离开的数据
包，它们应该被服务的相对顺序不会随着任何未来数据包的到达而改变。例如，在WFQ中就是这样，因
为任何数据包的QPS虚拟完成时间在数据包到达时就已经确定。然而，并非所有工作守恒算法都有这个
不变属性。例如，尽管WF2Q算法是工作守恒的（见第18节中的练习题10），但它没有这个不变属性
（见第18节中的练习题11）。
这个不变属性很重要，因为它与可编程性相关，因为具有这个属性的数据包调度算法可以使用一种称为
推入式先出（PIFO）队列的抽象数据结构来实现（Sivaraman等人，2016a）。PIFO是一个优先级队
列，允许元素根据元素的等级被推到任意相对位置，但总是从头部出队。在数据包调度使用工作守恒算
法的上下文中，这个等级是一个时间戳，指示数据包相对于其他数据包的优先级。例如，在WFQ中，数
据包的时间戳是其GPS虚拟完成时间。因此，在这个上下文中，PIFO实际上就是一个标准的优先级队
列。然而，使用PIFO这个术语（而不是简单地称之为优先级队列）有三个原因。首先，在数据包调度的
上下文中，PIFO不仅仅是一个抽象数据结构：它还体现了上述不变属性，即现有数据包的服务顺序不会
随着未来数据包的到达而改变。其次，正如Sivaraman等人（2016a）所解释的，对于具有这个不变属性
的非工作守恒算法，PIFO是一个日历队列。第三，PIFO这个术语最早是在Chuang等人（1999）中引入
的，用于描述所提出的组合输入和输出队列（CIOQ）交换算法（在第13.17节中描述），其中PIFO仅指
这个不变属性（而不是数据结构）。在后续内容中，我们将工作守恒的数据包调度算法称为PIFO兼容，
如果它具有这个属性。
对于每个支持的PIFO兼容算法，只有一个参数需要在可编程数据包调度框架下指定：一个回调过程，在
数据包到达时立即为每个数据包分配一个时间戳。在Sivaraman等人（2016b）中，这样的回调过程被称
为调度事务。例如，WFQ中的调度事务可以使用第12节中描述的有效的GPS时钟跟踪算法来实现。另一
个例子是，FIFO（先进先出）中的调度事务是微不足道的：数据包的时间戳仅仅是数据包的到达时间。
Sivaraman等人（2016a）展示了这个可编程数据包调度框架支持，除了WFQ和FIFO之外，还有其他几
种著名的算法（包括那些非工作守恒的），如令牌桶过滤（在第4节中描述）、分层数据包公平队列（即

                                                    37
WF2Q+）（Bennett和Zhang，1996a）、最少松弛时间优先（Leung，1989）、速率控制服务纪律
（Zhang和Ferrari，1994）和细粒度优先级调度（例如，最短作业优先）。
15.2 通用数据包调度器
可编程数据包调度意味着，当交换机的数据包调度算法因新的带宽分配标准而改变时，交换机必须重新
编程。然而，频繁地重新编程交换机会干扰生产网络的高速操作。Universal Packet Scheduler论文
（Mittal等人，2015）研究了公平性标准的改变是否必然需要交换机重新编程。是否存在一个通用数据包
调度器S，可以模仿（称为重放）任何数据包调度算法的带宽分配行为？如果这个问题的答案是肯定的，
那么我们只需要一次性地在可编程交换机上编程S，之后永远不需要重新编程它。
为了简要介绍通用数据包调度器的概念，我们必须引入一些术语和符号。考虑一个由路由器通过链路连
接的路由器网络，边界由入口和出口路由器组成。每个路由器使用某种数据包调度算法来调度其输出链
路上的数据包，不同的路由器可能使用不同的数据包调度算法。
我们用 {Aα } 表示这些路由器使用的数据包调度算法的集合，意思是路由器 α 使用算法 Aα 。考虑
一个数据包到达实例 {(p, i(p), path(p))∣p ∈ P } 到入口路由器，其中 P 是一组数据包， i(p) 是数
          ​                                                          ​




据包 p 到达入口路由器的时间， path(p) 是数据包 p 通过网络到出口路由器的路径。设 o(p) 是数
据包 p 的离开时间（从网络中），或者说数据包 p 离开相应的出口路由器的时间。那么
 {(path(p), i(p), o(p))∣p ∈ P } 被称为一个调度。
假设我们现在在这些路由器上使用另一组数据包调度算法 {A′α } ，给定相同的数据包到达实例的结果调
度是 {(path(p), i(p), o′ (p))∣p ∈ P } ，其中 o′ (p) 是数据包 p 在 {A′α } 下的离开时间。在Mittal
                                              ​




等人（2015）中，如果对于所有 p ∈ P ， o′ (p) ≤ o(p) ，则称 {A′α } 重放 {Aα } 。通用数据包调
                                                         ​




度器是这样的调度器，当被所有路由器使用时，可以被“操纵”来重放任何可行的调度，定义为可以由某
                                                     ​         ​




个 {Aα } 产生的调度。这里，用于“操纵”通用数据包调度器的信息限于通常用于数据包调度的头部值，
如优先级值或时间戳。例如，不允许数据包 p 在头部有一个“完整路径调度”，指定数据包 p 在路径
      ​




 (p) 上每个跳点的离开截止时间。
Mittal等人（2015）展示了以下理论结果。首先，在一些温和且常识性的假设下，一般情况下不存在通用
数据包调度器。其次，如果对于每个数据包 p ∈ P ，在其路径上最多有两个拥塞点，那么可以操纵
Least Slack Time First (LSTF) 算法（Leung, 1989）成为一个通用数据包调度器，其中操纵是在入口点
为其数据包头部分配适当的松弛（在调度中）值。拥塞点是数据包必须等待其传输机会的节点（路由
器）。
第三，如果每个数据包的路径上最多有一个拥塞点，那么简单的优先级调度算法也可以被操纵成一个通
用数据包调度器，其中操纵是为每个数据包 p 分配一个适当的优先级值。最后，如果一个数据包的路径
上可能有三个以上的拥塞点，那么同样不存在通用数据包调度器。第二和第三个结果在实践中可能是有
用的，因为在数据中心和广域网中，通常每个路径上最多有两个拥塞点：入口路由器和出口路由器。

                                                                           38
16 可扩展的公平队列
使用多个队列为每个流提供服务，我们已经看到：（1）一个常数时间算法（DRR）可以提供带宽保证，
即使是使用软件在QoS上，以及（2）一个对数时间开销算法可以提供带宽和延迟保证；此外，可以使用
第12节中描述的高级数据结构将对数开销降到最低。因此，在从小型边缘路由器到更大的骨干（核心）
路由器的网络设备中实现QoS似乎很容易。
不幸的是，即使是Thompson等人（1997年）对骨干路由器的较早研究也显示，大约有250,000个并发
流。随着流量的增加，我们预计这个数字将增长到一百万甚至更大，因为互联网的速度和流量都在增
加。在骨干路由器中为百万流保持状态可能是一项艰巨的任务。如果状态保存在SRAM中，所需的内存量
可能非常昂贵；如果状态保存在DRAM中，状态查找可能会很慢。
更有说服力的是，互联网扩展和聚合的支持者指出，目前的互联网路由仅使用大约1,000,000个前缀来覆
盖超过十亿个节点。为什么QoS需要如此多的状态，而IP的其他组件却不需要？特别是，虽然QoS状态可
能在今天是可管理的，但它可能对互联网的可扩展性构成严重威胁。就像前缀聚合多个IP地址的路由一
样，有没有办法聚合流状态？
聚合意味着骨干路由器将以相同的方式对待流组。聚合要求：（1）合理地将组成员视为相同对待，以及
（2）从数据包头部到聚合组有一个有效的映射。例如，在IP路由的情况下：（1）一个前缀聚合了许多
共享相同输出链路的地址，通常是因为它们位于相对地理区域内，以及（2）最长匹配前缀提供了一个从
头部目标地址到适当前缀的有效映射。
我们描述了三种提供聚合QoS的有趣提议：随机聚合（随机公平队列）；网络边缘的聚合（DiffServ）；
以及网络边缘的聚合与对行为不当的流的有效的监管（核心无状态公平队列）。
16.1 随机聚合
SFQ（McKenney，1991）背后的想法是，通过交易公平性的确定性来减少状态。在这个提议中，骨干路
由器保持一组固定的流队列，比如说125,000个，它们在这些队列上执行DRR。当数据包到达时，一些数
据包字段集（比如说，TCP和UDP流量的目标、源以及目标和源端口）被哈希到一个流队列。因此，假
设一个流是由用于哈希的字段集定义的，给定的流将总是被哈希到同一个流队列。因此，对于250,000个
并发流和125,000个流队列，大约2个流将共享同一个流队列或哈希桶。
随机公平队列有两个缺点。首先，不同的骨干路由器可以将流哈希到不同的组，因为路由器需要能够在
哈希分布不均匀时更改其哈希函数。其次，SFQ不允许某些流被区别对待（无论是在一个路由器内还是
在路由器之间），这是QoS的一个关键特性。因此，SFQ只提供某种可扩展的和统一的带宽公平性。
16.2 边缘聚合

                                                   39
DiffServ提议（Blake等人，1998）背后的三个想法是：通过将流聚合成类来放宽系统要求（P3），以减
少区分流的能力为代价；将类到核心路由器的映射从核心路由器转移到边缘路由器（P3c，空间上转移计
算）；以及在IP头部（P10）中传递类信息到核心路由器。
因此，边缘路由器将流聚合成类，并使用IP TOS字段中的标准化值标记数据包类。IP服务类型（TOS）
字段原本打算用于这样的用途，但它从未被标准化；像思科这样的供应商在其网络中使用它来表示流量
类，如IP语音，但没有标准化的流量类定义。DiffServ组概括并标准化了这种供应商行为，为正在标准化
的类保留了值。正在讨论的一个类是所谓的加速服务，其中为该类保留了一定的带宽。另一个是保证服
务，它在输出队列中为RED提供了较低的丢弃概率。
然而，关键是，在DiffServ中，骨干路由器的工作要轻松得多。首先，它们根据TOS字段中的少数字段值
将流映射到类。其次，骨干路由器只需要管理少量的队列，通常是每个类一个队列，有时是一个类内的
每个子类一个队列；例如，保证服务目前指定了类内的三个服务级别。然而，边缘路由器必须根据类似
于ACL的规则和可能检查整个头部来将流映射到类。不过，这是一个很好的权衡，因为边缘路由器的工
作速度较慢。
16.3 边缘聚合与监管
使用边缘聚合，两个流（比如说F1和F2）如果预留了带宽（比如说B1和B2），可能会被聚合成一个类，
该类名义上预留了一些带宽，该带宽是B >= B1 + B2，适用于类中的所有流。考虑图24。假设F1决定超
额订阅并以高于B的速率发送。图24中的边缘路由器ER可能目前有足够的带宽让F1和F2的所有数据包通
过。不幸的是，当这个聚合类到达骨干（核心）路由器CR时，假设核心路由器带宽受限，必须丢弃数据
包。
理想情况下，CR应该只丢弃超额订阅的流，如F1，并让F2的所有数据包通过。




FIGURE 24 If flows F1 and F2 are aggregated by the time they reach the core router CR, how can
the core router realize that F1 is oversubscribing without keeping state for each (unaggregated)
flow?



                                                                                              40
然而，CR如何判断哪些流是超额订阅的呢？它可以通过为所有通过的流保持状态来实现，但这将破坏可
扩展性。一个聪明的想法，称为核心无状态公平队列（Stoica等人，1998），观察到边缘路由器ER有足
够的信息来区分超额订阅的流。因此，ER可以使用原则P10，将信息传递到CR的数据包头部。
然而，CR应该如何处理超额订阅的流呢？丢弃所有这样标记的数据包可能过于严厉。如果有足够的带宽
供一些超额订阅的流使用，那么CR合理地丢弃与超额订阅程度成比例的数据包似乎是合理的。因此，ER
应该在数据包头部传递一个与流超额订阅程度成比例的值。为了实现这个想法，CR可以随机丢弃
（P3a），丢弃概率与超额订阅程度成比例。
虽然这有一些误差概率，但已经足够接近。最重要的是，随机丢弃可以在不保留任何流状态的情况下在
CR上实现。实际上，CR正在实现RED，但丢弃概率是基于边缘路由器设置的包头部字段计算的。
虽然核心无状态是一个不错的想法，但我们注意到，与SFQ（可以独立于路由器合作实现）和DiffServ
（已经获得了对其标准化的TOS字段使用的足够支持）不同，核心无状态公平队列目前只是一个研究提
案（Stoica等人，1998）。

17 总结
在本章中，我们解决了路由器的另一个主要实现瓶颈：调度数据包以减少拥塞的影响，并为某些流提供
公平性和QoS保证。我们从提供拥塞反馈的方案（如RED）开始，逐步上升到提供带宽和延迟QoS保证
的方案。我们还研究了如何使用DiffServ等技术在核心路由器上扩展QoS状态。
一个实际的路由器通常必须选择这些单独方案的组合。如今，许多路由器提供RED、近似公平丢弃
（AFD）、令牌桶监管和多个队列及DRR。然而，重点是，除了提供延迟保证的方案外，所有这些方案
都可以高效实现。
即使是提供出色延迟保证的方案也可以相当高效地实现。例如，我们已经展示了WFQ和WF2Q都可以通
过第12节中描述的高级形状数据结构，以仅O(logN)的时间复杂度（每个数据包）提供出色的延迟保证。
此外，第14节中描述的快速公平队列（QFQ）可以以O(1)的时间复杂度提供类似的延迟保证。许多组合
方案也可以使用我们概述的原则高效实现。练习探讨了其中一些组合。
为了使本章自成一体，我们花了大量篇幅讨论诸如拥塞控制和资源预留等主题，这些主题实际上与本书
的主要内容是外围的。我们真正关心的是使用我们的原则来解决调度瓶颈。
为了避免忘记这一点，我们一如既往地提醒您，表1总结了本章中使用的技术和相应的原则。




                                                 41
5 通用处理器共享
通用处理器共享（GPS）是一种理想化的调度算法，属于公平排队（Fair Queuing）算法家族的理
论基础。GPS为多任务系统提供了一种完美的公平资源共享模型，其核心思想是通过虚拟时间机
制，为每个数据流提供独立、隔离的虚拟服务器，确保所有活动流能够同时获得公平的带宽分配。
GPS模型由Parekh和Gallager在1993年提出，已成为服务质量（QoS）保障的理论基准。
在 GPS中，一个调度器处理 N 个流（也称为“classes”或“sessions”），并为每个流配置一个权重
wi 。然后，GPS 确保对于任意一个流 i 和某个时间区间 (s, t] ，如果流 i 在此区间内持续处于
飞空状态（即队列从未为空），那么对于任何其他流 j ，以下关系成立：
  ​




                      wj Oi (s, t) ≥ wi Oj (s, t)
                         ​    ​                   ​           ​




其中， Ok (s, t) 表示流 k 在时间区间 (s, t] 内输出的比特数。
可以证明，每个流 i 至少获得如下速率：
      ​




                                        wi
                             Ri =             R
                                          ​




                                       N
                                  ​                       ​




                                      ∑j=1 wj ​       ​




其中， R 是服务器的总速率。
这是一个最小速率。如果在某一时间段内，某些流未使用其带宽，则剩余的容量将由活动流根据其各自
的权重进行分配。例如，考虑一个 GPS 服务器，其权重设置为 w1 = 2, w2 = w3 = 1 。在这种情况
下，第一个流将至少获得一半的容量，而另外两个流各获得 1/4 的容量。然而，如果在某个时间区间
                                                                  ​   ​   ​




(s, t] 内，只有第二个和第三个流是活动的，它们将各自获得一半的容量。

5.1 一个简单的例子
以下通过一个 网络带宽分配的例子 说明通用处理器共享（GPS）的工作机制，帮助理解其核心原理和公
平性逻辑。
场景设定
假设服务器总带宽为 ( r = 100 Mbps})，需要调度 3个数据流（流A、流B、流C），每个流分配的权重分
别为：
 ● 流A：( wA = 5 )
          ​




                                                                              42
 ● 流B：( wB = 3 )            ​




 ● 流C：( wC = 2 )                ​




总权重为 ( ∑ wi = 5 + 3 + 2 = 10 )。     ​




根据GPS模型，每个流的最小速率为： ri = ∑wwi i ⋅ r              ​
                                                    ​




                                                            ​




计算得：
                                                        ​




        5
 ● 流A：    × 100 = 50, Mbps
       10
                ​




 ● 流B：
        3
          × 100 = 30, Mbps
       10
                    ​




 ● 流C：
        2
          × 100 = 20, Mbps
       10
                        ​




情况1：所有流持续积压（满负荷传输）
假设三个流始终有数据需要传输（队列不为空），则它们严格按照权重分配带宽，实际速率等于最小速
率：
 ● 流A：50 Mbps（占50%）
 ● 流B：30 Mbps（占30%）
 ● 流C：20 Mbps（占20%）



情况2：某一流空闲（部分流未占满带宽）
假设流B在某段时间内空闲（无数据传输），则剩余带宽由活动流（流A和流C）按权重重新分配。
此时总活动权重为 wA + wC = 5 + 2 = 7 ，新的速率为：    ​   ​




 ● 流A：
       5
         × 100 ≈ 71.43, Mbps
       7
        ​




 ● 流C：
       2
         × 100 ≈ 28.57, Mbps
       7
            ​




 ● 流B传输数据：0 Mbit
 ● 总传输量：约100 Mbit（仍占满服务器带宽，流A和流C按权重“瓜分”流B的空闲带宽）




6.权重公平队列（WFQ）
                                                                43
用于在多个数据流之间公平分配带宽资源。它通过为不同数据流设置 权重（Weight）来实现差异化服
务，既保证低权重流的基本带宽需求，又允许高权重流在资源空闲时占用更多带宽。WFQ 是 通用处理
器共享（GPS）模型的近似实现，适用于实际网络中变长数据包的场景。
1. 权重分配与带宽比例
 ● 权重（Weight）的定义
   每个队列 i 被赋予一个正实数权重 wi ，代表该队列对带宽资源的“优先级”或“资源占比权值”。
   权重越大，队列获得的带宽比例越高。
                                             ​




 ● 带宽分配公式

   假设总可用带宽为 R ，所有队列的权重之和为 ∑ wj ，则队列 i 的 最小保证带宽 为：​




               wi
     ri =          ⋅R
                ​




              ∑ wj
         ​              ​




                    ​




     ○       动态共享机制：当某队列无数据发送时，其带宽会按权重比例分配给其他活跃队列。例如，若
             队列B空闲，队列A可独占全部 (100Mbps)。
2. 虚拟完成时间（Virtual Finish Time, VFT）
WFQ 通过 虚拟完成时间 VFT 将离散数据包的处理顺序映射到 GPS 的理论完成时间，证明在数据包到
达服从一定条件（如独立同分布）时，WFQ 的队列延迟上界与 GPS 一致，从而实现近似公平性。
VFT的计算规则
 ● 若队列 i 为空，新数据包的 VFT = 当前时间 + 数据包大小/ wi 。                           ​




 ● 若队列 i 非空，新数据包的 VFT = 前一个数据包的 VFT + 数据包大小/ wi 。                        ​




调度器始终选择 VFT 最小的数据包 优先发送，确保高权重流的数据包（或小尺寸数据包）优先传
输，同时保证低权重流不会被饿死。

设当前系统时间为 t ，处理数据包的顺序为其进入队列的顺序，每个数据包 p 的长度为 Lp ，所属队
列为 i ,权重为 wi 。
                                                                             ​




                            ​




 ● 情况1：队列 i 为空
   新数据包的虚拟完成时间为：
                            Lp
     VFT = t +
                                     ​




                                         ​




                            wi   ​




     ○       物理意义：假设从当前时间 t 开始，持续占用 ∑wwi ⋅ R 带宽处理该数据包，其完成时间
                                                     ​




                                                                 ​




                                                         j   ​




                                                                                 44
         为 t + Lrip ，而 ri = ∑wwi j ⋅ R ，因此 Lrip = Lpw⋅i∑⋅ Rwj 。
                 ​
                     ​




                         ​           ​
                                                          ​




                                                              ​




                                                                                                  ​
                                                                                                          ​




                                                                                                              ​
                                                                                                                                      ​




                                                                                                                                          ​
                                                                                                                                                  ​




                                                                                                                                                      ​




     为简化计算，WFQ 直接使用 Lwpi （忽略总带宽 R ），最终调度顺序与理论模型一致。
     ○
                                                                      ​




                                                                          ​




 ● 情况2：队列 i 非空
                                                                  ​




   假设队列中前一个数据包的虚拟完成时间为 VFTprev ，则新数据包的虚拟完成时间为：
                                     Lp
     VFT = VFTprev +
                                              ​




                                     wi
                                                  ​




         物理意义：模拟前一个数据包处理完毕后，立即开始处理当前数据包的“流体”过程。
                                          ​




     ○


3. 调度策略：基于VFT的优先级队列
 ● 全局优先级队列
   WFQ 维护一个 全局队列，存储所有队列中待发送的数据包，并按 虚拟完成时间（VFT）升序排
   列。
 ● 调度规则
    ○ 每次从全局队列中选取 VFT 最小的数据包 发送，确保“理论上应先完成的数据包”优先传输。
    ○ 若多个数据包的 VFT 相同，则按队列优先级或到达顺序（如 FIFO）处理。


示例：调度过程解析
假设现有两个队列：
 ● 队列1（权重 w1 = 2 ：数据包A（长度4）、数据包C（长度6），依次到达。
                             ​




 ● 队列2（权重 w2 = 1 ：数据包B（长度3），在数据包A之后到达。
                                 ​




计算各数据包的VFT：
 1. 数据包A（队列1，第一个数据包） VFTA = t0 +
                                     4
                                       = t0 + 2
                                     2
                                                                                      ​       ​                       ​           ​




                                       = t1 + 3 (假设 t1 = t0 + 1)
                                    3
 2. 数据包B（队列2，第一个数据包） VFTB = t1 +
                                    1
                                                                                  ​       ​                       ​           ​                           ​       ​




 3. 数据包C（队列1，第二个数据包） VFTC = VFTA +
                                        6
                                           = (t0 + 2) + 3 = t0 + 5
                                        2
                                                                              ​                       ​                   ​                   ​               ​




调度顺序：比较 VFT 值，假设 t0 = 0 ，则：                           ​




 ● 数据包A（VFT=2）→ 数据包B（VFT=4）→ 数据包C（VFT=5）。
 ● 实际传输顺序与 GPS 模型中“按权重比例连续分配资源”的理论结果一致：
     ○ 队列1在 t = 0 ∼ 2 处理数据包A（占用 2/3 带宽），
     ○ 队列2在 t = 2 ∼ 5 处理数据包B（占用 1/3 带宽，需时间 3/(1/3R) = 9/R ，假设



                                                                                                                                                                      45
           R=1         ，则耗时3单位时间，与 VFT 一致）。

4. 公平性与饥饿避免
 ●   公平性的数学本质
     WFQ 通过 VFT 确保任意时刻 t ，队列 i 已获得的服务量（即已传输的数据包总长度）满足：
     Si (t)    Sj (t)
            ≤                      (∀i, j)
       ​               ​




      wi         wj
               ​               ​




   其中     为队列 i 在时间 t 内的累计服务量。这保证了各队列的服务量与权重成正比，实
           ​               ​




          Si (t)
   现“按比例公平”（Proportional Fairness）。
                   ​




 ● 饥饿避免
   由于每个队列的数据包都会被分配有限的 VFT（只要队列持续有数据到达），低权重队列的数据包
   不会因高权重队列的持续流量而无限期延迟，从而避免“饿死”。
5. 与GPS模型的关系
 ● GPS的理想假设
   假设数据以“流体”形式连续到达和处理，每个队列 (i) 的数据被按权重比例 (w_i / \sum w_j) 实时
   分配带宽，无分组开销。
 ● WFQ的近似性
   通过 VFT 将离散数据包的处理顺序映射到 GPS 的理论完成时间，证明在数据包到达服从一定条件
   （如独立同分布）时，WFQ 的队列延迟上界与 GPS 一致，从而实现近似公平性。

WFQ 的问题在于，在某些情况下，它无法很好地近似 GPS（广义处理器共享）。在以下示例中，所有
数据包的大小均为 1，数据速率是每时间单位 1 个数据包。共有六个流，其权重如下：流 F1 被分配了总
输出链路的 50% 权重，而流 F2 至 F6 各自被分配了 10% 的权重。此时，F1 流同时到达了六个数据
包，而 F2 至 F6 每个流各到达了一个数据包。




                                                           46
GPS 在 t=0 到 t=10 时间段内处理来自 F2 至 F6 的数据包。此外，它还在 t=0 到 t=2 处理了 F1 的第
一个数据包（A），接着从 t=2 到 t=4 处理数据包 B，以此类推。




WFQ 以图中所示的方式处理这些数据包。注意，在处理完 E 数据包后出现了一个空隙。WF2Q（或称
Worst Case Fair Weighted Queue）也被展示出来，并且在高权重数据包之间没有大的间隔。
WF2Q （Worst Case Fair Fair Weighted Queue）通过对 WFQ 算法进行轻微修改实现了这一点。与
WFQ 类似，WF2Q 也使用了虚拟时间的概念。虚拟完成时间是指 GPS 完成发送该数据包的时间。然
而，与 WFQ 不同的是，WF2Q 并不是寻找队列中所有数据包的最小虚拟完成时间，而是寻找那些虚拟
开始时间已经过去的、具有最小虚拟完成时间的数据包。虚拟开始时间是指 GPS 开始发送该数据包的时
间。虚拟开始时间和虚拟完成时间定义如下：
F(i, k, t) = MAX{F(i, k-1, t), R(t)} + P(i, k, t)
S(i, k, t) = MAX{F(i, k-1, t), R(t)}
其中：
  ● i = 连接编号


                                                                47
 ●  k = 数据包编号
 ● t = 时间
 ● R(t) = 轮次编号（虚拟时间）
 ● P(i, k, t) = 在时间 t 到达的连接 i 上第 k 个数据包以连接 i 的带宽进行传输所需的时间
 ● F(i, k-1, t) = 连接 i 上第 (k-1) 个数据包的完成编号

解释：
 1. F(i, k, t) 表示在时间 t，连接 i 上第 k 个数据包的完成时间。它由两部分组成：
      ○ 前一个数据包的完成时间 F(i, k-1, t) 和当前轮次编号 R(t) 中的最大值，表示数据包的起始时
         间；
      ○ 再加上当前数据包的传输时间 P(i, k, t)。

 2. S(i, k, t) 表示在时间 t，连接 i 上第 k 个数据包的开始时间。它等于前一个数据包的完成时间 F(i,
    k-1, t) 和当前轮次编号 R(t) 的最大值。
 3. R(t) 是虚拟时间的概念，通常用于公平队列调度算法中，表示当前系统的虚拟进度。

 4. P(i, k, t) 取决于连接 i 的带宽以及第 k 个数据包的大小，用来计算该数据包的传输时间。

 5. F(i, k-1, t) 表示连接 i 上前一个数据包的完成时间，它是计算当前数据包完成时间的基础。




                                                                 48
